\chapter{Primes on Average}

\begin{center}\itshape
[The primes] grow like weeds among the natural numbers, seeming to obey no other law than that of chance, and nobody can predict where the next one will sprout.
[\dots] The prime numbers exhibit stunning regularity, [\dots] there are laws governing their behaviour, and [\dots] they obey these laws with almost military precision.\\
\upshape --- Don Zagier, \emph{The first 50 million prime numbers}, inaugural lecture at Bonn (1975), translated by R.~Perlis
\end{center}

One day Kannan Soundararajan, of Stanford University, went to a lecture by Tadashi Tokieda about tossing coins. Alice tosses a fair coin until she sees heads followed by
tails; Bob tosses until he sees two heads in a row. Both patterns have probability $\frac14$ on any two given tosses, and yet Alice waits four tosses on average and Bob six.
A pattern that overlaps itself, like heads-heads, arrives in clumps, and clumps make the waits longer. Soundararajan wondered whether the primes hide quirks of the same kind,
and with Robert Lemke Oliver he looked at the last digits of consecutive primes. Apart from $2$ and $5$, every prime ends in $1$, $3$, $7$ or $9$, and by the prime
number theorem for progressions the four endings are equally frequent in the long run. One would expect the ending of the \emph{next} prime to be just as indifferent to the ending of the last one.
It is not. Among the primes below $10^8$, a prime ending in $1$ is followed by another prime ending in $1$ only $17.7\%$ of the time, and by a prime ending in $7$
$31.0\%$ of the time. In 2016 Lemke Oliver and Soundararajan explained the size of the effect with the conjectures of Hardy and Littlewood. ``The primes really hate
to repeat themselves'', Lemke Oliver told \emph{Quanta Magazine}. Problem~\ref{pb:digitpairs} builds a coin-tossing model that accounts for most of the bias, and shows
what it misses.

The story has both halves of Zagier's sentence in it. Taken one at a time, the primes look like weeds; taken on average, they obey laws with military precision, and the
deviations from the averages are themselves lawful. This chapter is a walk among those averages. We add up primes and prime factors, look at the numbers $p-1$ and
$n^2+1$ through the eyes of the prime number theorem, meet the smooth and the rough numbers and the strange constant $e^{-\gamma}$ that haunts them, and end with
distributions, gaps and digits. The tools are the prime number theorem and its version in arithmetic progressions (Volume~5, Theorems~5.5 and~5.8), the three theorems of
Mertens (Volume~5, Chapter~1), the sieve bounds of Volume~5, Chapter~8, and the probability and analysis of Volumes~3 and~4.

\begin{think}
Below $10^8$, after a prime ending in $1$, guess how often the next prime ends in $1$, $3$, $7$ and $9$. Then a second guess: for a random integer $n\le x$, is
the average size of its largest prime factor of order $\frac x2$, of order $\frac x{\log x}$, or of order $\sqrt x$?
\end{think}

Throughout, $p$ and $q$ denote primes, $\pi(x)$ is the number of primes up to $x$, $\pi(x;d,a)$ the number of those that are $\equiv a\pmod d$, $p_n$ is the $n$-th
prime, $\omega(n)$ and $\Omega(n)$ count the prime factors of $n$ without and with multiplicity, $P^+(n)$ and $P^-(n)$ are the largest and smallest prime factors of
$n\ge2$, and $\gamma=0.5772\ldots$ is Euler's constant. Two consequences of Volume~5 are used again and again. Chebyshev's bound $\pi(y)\le C\frac y{\log y}$ for $y\ge2$
(Volume~5, Proposition~1.2), and the Brun--Titchmarsh inequality in progressions (Volume~5, Exercise~8.4), in the form: there is $x_0$ such that
\begin{equation}\label{eq:bt13}
\pi(x;d,1)\le\frac{5x}{\varphi(d)\log x}\quad(x\ge x_0,\ 1\le d\le\sqrt x).
\end{equation}

\section{Adding up primes and prime factors}

A sum over primes is an integral against $\mathrm d\pi(t)$, and the prime number theorem turns it into an ordinary integral. The first problems apply this to sums of
primes and of prime factors, and then to the prime factors of two famous products, $n!$ and $\binom{2n}n$.

\begin{exercise}[$\star$]\label{pb:leastpf}
\emph{Sums of primes, and of least prime factors.}
\begin{enumerate}[label=(\alph*)]
\item Show that $\sum_{p\le x}p\sim\frac{x^2}{2\log x}$.
\item Show that $\sum_{2\le n\le x}P^-(n)\sim\frac{x^2}{2\log x}$: the least prime factors of all the integers up to $x$ add up to essentially the sum of the primes alone.
\item Show that the sum of the first $n$ primes is asymptotic to $\frac12n^2\log n$.
\item Numerically, one finds $\sum_{p\le10^7}p=3\,203\,324\,994\,356$, which is $1.0326\,\frac{x^2}{2\log x}$ but $0.99950\operatorname{Li}(x^2)$ for $x=10^7$. Explain why
$\operatorname{Li}(x^2)$ is the natural approximation, and show that $\sum_{p\le x}p=\operatorname{Li}(x^2)+O(x^2e^{-c\sqrt{\log x}})$ for some $c>0$.
\end{enumerate}
\end{exercise}

\begin{exercise}[$\star\star\star$]\label{pb:largestpf}
\emph{The largest prime factor on average.} For $n\ge2$ let $A(n)$ be the sum of the prime factors of $n$ counted with multiplicity, so $A(12)=2+2+3=7$. For $y\ge1$ and
$z\ge2$, let $\Psi(y,z)$ be the number of $m\le y$ with no prime factor larger than $z$ (with $m=1$ counted).
\begin{enumerate}[label=(\alph*)]
\item Show that $\sum_{2\le n\le x}P^+(n)=\sum_{p\le x}p\,\Psi(\frac xp,p)$.
\item Show that the primes $p\le\sqrt x$ contribute $O(x^{3/2})$, and that $\Psi(\frac xp,p)=\lfloor\frac xp\rfloor$ when $p>\sqrt x$.
\item Show that $\sum_{\sqrt x<p\le x}p\lfloor\frac xp\rfloor=\sum_{1\le k<\sqrt x}\Bigl(\sum_{p\le x/k}p-\sum_{p\le\sqrt x}p\Bigr)$, and deduce that
\[
\sum_{2\le n\le x}P^+(n)\sim\frac{\pi^2}{12}\,\frac{x^2}{\log x}.
\]
\item Show that $A(n)-P^+(n)\le\sqrt n\log_2n$, and deduce the same asymptotic formula for $\sum_{n\le x}A(n)$.
\item Numerically, one finds $\sum_{n\le10^7}P^+(n)=1.0767\cdot\frac{\pi^2}{12}\frac{x^2}{\log x}$. Why is the ratio larger than $1$, and why does it approach $1$ so slowly?
\end{enumerate}
\end{exercise}

\begin{exercise}[$\star\star$]\label{pb:radical}
\emph{The average radical.} The radical $\operatorname{rad}(n)$ is the product of the distinct primes dividing $n$, so $\operatorname{rad}(72)=6$. It is the quantity in
the $abc$ conjecture of Masser and Oesterl\'e. Let $C=\prod_p\bigl(1-\frac1{p(p+1)}\bigr)=0.70444\ldots$
\begin{enumerate}[label=(\alph*)]
\item Write $\frac{\operatorname{rad}(n)}n=\sum_{d\mid n}g(d)$ and compute $g$. Show that $g(d)=0$ unless $d$ is \emph{squarefull} (every prime factor of $d$ appears at
least squared), and that $|g(d)|\le1$.
\item Show that the number of squarefull $d\le y$ is at most $\zeta(\frac32)\sqrt y$.
\item Deduce that $\sum_{n\le x}\operatorname{rad}(n)=\frac C2x^2+O(x^{3/2})$, and that the mean value of $\frac{\operatorname{rad}(n)}n$ is $C$.
\end{enumerate}
\end{exercise}

\begin{exercise}[$\star\star$]\label{pb:factorialdiv}
\emph{How many divisors has $n!$?} Let $\tau(m)$ be the number of divisors of $m$, and $C=\sum_{k\ge1}\frac1k\log\bigl(1+\frac1k\bigr)=1.25774\ldots$
\begin{enumerate}[label=(\alph*)]
\item Show that $\tau(n!)=\prod_{p\le n}\bigl(1+v_p(n!)\bigr)$, where $v_p(n!)=\sum_{j\ge1}\lfloor\frac n{p^j}\rfloor$.
\item Show that the primes $p\le\sqrt n$ contribute $O(\sqrt n)$ to $\log\tau(n!)$, and that $v_p(n!)=\lfloor\frac np\rfloor$ for $p>\sqrt n$.
\item Show that $\log\tau(n!)\sim C\frac n{\log n}$. In other words, $\tau(n!)=e^{C(1+o(1))n/\log n}$, where $e^C=2\cdot\sqrt{\frac32}\cdot\sqrt[3]{\frac43}\cdots=3.5174\ldots$
\item Numerically, one finds $\log\tau(n!)\big/\frac n{\log n}=1.58$, $1.51$, $1.45$, $1.42$ for $n=10^4,10^5,10^6,10^7$. Comment.
\end{enumerate}
\end{exercise}

\begin{exercise}[$\star\star$]\label{pb:centralbinom}
\emph{The prime factors of $\binom{2n}n$.} Every prime in $]n,2n]$ divides $\binom{2n}n$, and no prime in $]\frac{2n}3,n]$ does once $n\ge5$ (this is the heart of Erd\H os's proof of Bertrand's postulate).
\begin{enumerate}[label=(\alph*)]
\item Show that for $p>\sqrt{2n}$, $v_p\bigl(\binom{2n}n\bigr)=\lfloor\frac{2n}p\rfloor-2\lfloor\frac np\rfloor$, which is $1$ if the fractional part of $\frac np$ is at least $\frac12$, and $0$ otherwise.
\item Deduce that the primes $p>\sqrt{2n}$ dividing $\binom{2n}n$ are those of $\bigcup_{k\ge0}\bigl]\frac n{k+1},\frac n{k+1/2}\bigr]$.
\item Show that $\omega\bigl(\binom{2n}n\bigr)\sim\log4\cdot\frac n{\log n}$: the central binomial coefficient has about $39\%$ more prime factors than the primes of $]n,2n]$ alone.
\item Numerically, one finds $\omega\bigl(\binom{2n}n\bigr)=453\,587$ for $n=5\cdot10^6$. Compare.
\end{enumerate}
\end{exercise}

\begin{exercise}[$\star\star$]\label{pb:primezeta}
\emph{Computing $\sum_p\frac1{p^2}$ to twelve digits.} For real $s>1$ let $P(s)=\sum_pp^{-s}$.
\begin{enumerate}[label=(\alph*)]
\item Show that $\log\zeta(s)=\sum_{k\ge1}\frac{P(ks)}k$.
\item Show that $P(s)=\sum_{k\ge1}\frac{\mu(k)}k\log\zeta(ks)$, the double series converging absolutely.
\item Show that if the sum in (b) is stopped at $k=K$, the error is at most $\frac{4}{K}2^{-(K+1)s}$ for $s\ge2$.
\item Summing the primes up to $10^7$ gives $\sum_{p\le10^7}p^{-2}=0.4522474141\ldots$, wrong in the eighth decimal. With (b) and $K=20$, one finds numerically
$P(2)=0.452247420041\ldots$, correct to twelve digits. Estimate the tail $\sum_{p>10^7}p^{-2}$ and check the two numbers against each other.
\end{enumerate}
\end{exercise}

\begin{exercise}[$\star$]\label{pb:superprimes}
\emph{Primes with a prime index.} The primes $p_{p_n}$ are $3,5,11,17,31,41,59,\dots$
\begin{enumerate}[label=(\alph*)]
\item Show that $p_{p_n}\sim n\log^2n$ and that the number of them up to $x$, $\pi(\pi(x))$, is asymptotic to $\frac x{\log^2x}$.
\item Show that $\sum_n\frac1{p_{p_n}}$ converges. Up to $10^8$ the partial sum is $0.9833$; how large, roughly, is the remaining tail?
\item Numerically, one finds $\pi(\pi(10^8))=397\,557$, which is $1.349\,\frac x{\log^2x}$. Find a better approximation using $\pi(y)\approx\frac y{\log y}(1+\frac1{\log y})$.
\end{enumerate}
\end{exercise}

\begin{exercise}[$\star\star$]\label{pb:ramprimes}
\emph{Ramanujan primes.} In 1919 Ramanujan gave a short proof of Bertrand's postulate which shows more: $\pi(x)-\pi(\frac x2)\to\infty$. Let $f(x)=\pi(x)-\pi(\frac x2)$ for real $x>0$.
\begin{enumerate}[label=(\alph*)]
\item For $n\ge1$ let $R_n$ be the least real number such that $f(x)\ge n$ for all $x\ge R_n$. Show that $R_n$ exists and is a prime, and that $f(R_n)=n$.
(So $R_1=2$, $R_2=11$, $R_3=17$, $R_4=29$, $R_5=41$; Sondow named them Ramanujan primes in 2009.)
\item Show that $R_n\sim p_{2n}$.
\item Numerically, one finds $R_{10\,000}=242\,057$ while $p_{20\,000}=224\,737$. Using $\pi(y)\approx\frac y{\log y}(1+\frac1{\log y})$, show that one should expect
$\frac{R_n}{p_{2n}}\approx1+\frac{\log2}{\log p_{2n}}$, and compare.
\end{enumerate}
\end{exercise}

\section{Primes in disguise: \texorpdfstring{$p-1$}{p-1}, \texorpdfstring{$n^2+1$}{n\^{}2+1} and \texorpdfstring{$x^n-1$}{x\^{}n-1}}

Pick a random prime $p$ and look at $p-1$. It is an even number, but beyond that, how does it factor? Is it squarefree as often as a random integer, $61\%$ of the time?
How many divisors does it have? The answers need the primes in the progressions $1\bmod d$ for many $d$ at once, and the Brun--Titchmarsh inequality~\eqref{eq:bt13}
is exactly the tool that keeps the large $d$ under control.

\begin{exercise}[$\star\star\star$]\label{pb:pminusone}
\emph{Artin's constant twice.} Let $A=\prod_q\bigl(1-\frac1{q(q-1)}\bigr)=0.3739558\ldots$, the product over all primes $q$.
\begin{enumerate}[label=(\alph*)]
\item (\emph{The large squares are rare.}) Show that for $y\ge2$ and $x\ge x_0$, the number of primes $p\le x$ such that $q^2\mid p-1$ for some prime $q>y$ is at most $\frac{10x}{y\log x}+2x^{3/4}$.
\item Show that the proportion of primes $p\le x$ for which $p-1$ is squarefree tends to $A$.
\item Show that $\frac1{\pi(x)}\sum_{p\le x}\frac{\varphi(p-1)}{p-1}\to A$ as well.
\item Explain why (c) says that a random residue modulo a random prime is a primitive root with probability $A$. Numerically, one finds that $2$ is a primitive root modulo
$248\,491$ of the $664\,578$ odd primes below $10^7$, a proportion of $0.37391$. (That this proportion tends to $A$ is Artin's conjecture, still unproved.)
\end{enumerate}
\end{exercise}

\begin{exercise}[$\star\star$]\label{pb:titchmarsh}
\emph{The divisors of $p-1$.} Let $Z=\frac{\zeta(2)\zeta(3)}{\zeta(6)}=1.94359\ldots$ and $B=\prod_q\bigl(1+\frac1{(q-1)^2}\bigr)=2.82642\ldots$
\begin{enumerate}[label=(\alph*)]
\item Show that $\sum_{d\le y}\frac1{\varphi(d)}\le c(1+\log y)$ for some constant $c$, writing $\frac1{\varphi(d)}=\frac1d\sum_{e\mid d}\frac{\mu(e)^2}{\varphi(e)}$.
\item Show that $\sum_{p\le x}\tau(p-1)=O(x)$.
\item Show that $\frac1{\pi(x)}\sum_{p\le x}\frac{p-1}{\varphi(p-1)}\to B$.
\item Assuming that $\sum_{d\le y}\frac1{\varphi(d)}\sim Z\log y$ (a theorem of Landau), give a heuristic argument that $\sum_{p\le x}\tau(p-1)\sim Zx$. Numerically, one finds
$\frac1x\sum_{p\le x}\tau(p-1)=1.916$, $1.927$, $1.932$ for $x=10^5,10^6,10^7$.
\end{enumerate}
\end{exercise}

\begin{exercise}[$\star\star$]\label{pb:menon}
\emph{The roots of $x^n-1$ modulo $p$, on average.} Fix $n\ge1$.
\begin{enumerate}[label=(\alph*)]
\item Show that for every prime $p$, the polynomial $x^n-1$ has exactly $\gcd(n,p-1)$ roots in $\mathbb F_p$.
\item (\emph{Menon's identity, 1965.}) Let the group $(\Z/n\Z)^\times$ act on $\Z/n\Z$ by multiplication. Count the orbits in two ways and deduce that
\[
\sum_{a\in(\Z/n\Z)^\times}\gcd(a-1,n)=\varphi(n)\tau(n).
\]
\item Deduce that the average number of roots of $x^n-1$ modulo $p$, over the primes $p\le x$, tends to $\tau(n)$.
\item Show directly that for $d\mid n$, the cyclotomic polynomial $\Phi_d$ has on average one root modulo $p$, and explain (b) and (c) through the factorisation
$x^n-1=\prod_{d\mid n}\Phi_d(x)$. Numerically, one finds the averages $5.994$, $7.983$ and $11.979$ for $n=12,24,60$, over the primes below $10^6$.
\end{enumerate}
\end{exercise}

\begin{exercise}[$\star\star\star$]\label{pb:omegapoly}
\emph{How many prime factors has $n^2+1$?} Let $\rho(p)$ be the number of solutions of $n^2\equiv-1\pmod p$, and $\chi$ the non-principal character modulo $4$.
\begin{enumerate}[label=(\alph*)]
\item Show that $\rho(2)=1$ and $\rho(p)=1+\chi(p)$ for odd $p$.
\item Show that $n^2+1$ has at most one prime factor larger than $x$ when $n\le x$, and deduce that $\sum_{n\le x}\omega(n^2+1)=x\sum_{p\le x}\frac{\rho(p)}p+O(x)$.
\item Deduce that $\sum_{n\le x}\omega(n^2+1)=x\log\log x+O(x)$.
\item Show that $\sum_{n\le x}\bigl(\omega(n^2+1)-\log\log x\bigr)^2=O(x\log\log x)$, and deduce that $\omega(n^2+1)$ has normal order $\log\log n$.
\item Numerically, one finds $\frac1x\sum_{n\le x}\omega(n^2+1)-\log\log x=0.598$, $0.603$, $0.606$ for $x=10^4,10^5,10^6$. Is this consistent with (c)?
\end{enumerate}
\end{exercise}

\section{Smooth numbers, rough numbers and the constant \texorpdfstring{$e^{-\gamma}$}{e\^{}(-gamma)}}

An integer is \emph{$y$-smooth} if all its prime factors are at most $y$, and \emph{$y$-rough} if all of them exceed $y$. Mertens's third theorem (Volume~5, Theorem~1.6),
\[
\prod_{p\le y}\Bigl(1-\frac1p\Bigr)\sim\frac{e^{-\gamma}}{\log y},
\]
says how many integers survive a sieve by the primes up to $y$, \emph{if the survivors are counted in a long enough interval}. The trouble starts when the interval is
not long enough, and Euler's constant appears where nobody invited it. The next five problems chase it.

We write $\Psi(x,y)$ for the number of $y$-smooth $n\le x$ and $\Phi(x,y)$ for the number of $y$-rough $n\le x$, both counting $n=1$, and $u=\frac{\log x}{\log y}$.
\emph{Dickman's function} $\rho$ is the continuous function on $[0,\infty[$ with $\rho(u)=1$ for $0\le u\le1$ and $u\rho'(u)=-\rho(u-1)$ for $u>1$. It is the limit of
$\frac1x\Psi(x,x^{1/u})$, but we shall not need this.

\begin{exercise}[$\star\star$]\label{pb:rankin}
\emph{Rankin's trick.}
\begin{enumerate}[label=(\alph*)]
\item Show that $\Psi(x,y)\le x^\sigma\prod_{p\le y}(1-p^{-\sigma})^{-1}$ for every $\sigma\in\left]0,1\right]$.
\item Take $\sigma=1-\frac\eta{\log y}$ with $0<\eta\le\frac13\log y$. Using $e^{\eta t}\le1+(e^\eta-1)t$ for $0\le t\le1$ and the first two theorems of Mertens, show that
there is an absolute constant $C$ with
\[
\Psi(x,y)\le Cx\log y\,\exp\bigl(-\eta u+Ce^\eta\bigr).
\]
\item Choose $\eta=\log u$ and deduce that $\Psi(x,y)\le Cx\log y\,e^{-u\log u+Cu}$ when $u\ge1$ and $\log u\le\frac13\log y$. In particular, for $y=e^{\sqrt{\log x}}$,
the proportion of $y$-smooth integers up to $x$ is at most $e^{-(\frac12+o(1))\sqrt{\log x}\log\log x}$. (This is the scale on which the quadratic sieve lives.)
\item Direct computation gives $\Psi(10^{10},100)=8\,800\,084$ and $\Psi(10^{12},1000)=6\,471\,274\,933$, and minimises the bound of (a) over $\sigma$: it gets
$9.95\cdot10^7$ and $1.26\cdot10^{11}$. By what factor does the best Rankin bound miss? Where is the loss?
\end{enumerate}
\end{exercise}

\begin{exercise}[$\star\star\star$]\label{pb:dickmanint}
\emph{Dickman's function integrates to $e^\gamma$.}
\begin{enumerate}[label=(\alph*)]
\item Show that $u\rho(u)=\int_{u-1}^u\rho(t)\dd t$ for $u\ge1$. Deduce that $\rho>0$, that $\rho$ decreases on $[1,\infty[$, and that $\rho(u)\le\frac12\rho(u-1)$ for $u\ge2$.
\item Let $\hat\rho(s)=\int_0^\infty e^{-su}\rho(u)\dd u$ for $s\ge0$. Extending $\rho$ by $0$ on $]-\infty,0[$, show that $(u\rho(u))'=\rho(u)-\rho(u-1)$ for $u>0$, $u\ne1$, and deduce
that $s\hat\rho'(s)=-(1-e^{-s})\hat\rho(s)$ for $s>0$.
\item Deduce that $\hat\rho(s)=\hat\rho(0)\exp\bigl(-\int_0^s\frac{1-e^{-t}}t\dd t\bigr)$.
\item Show that $s\hat\rho(s)\to1$ as $s\to\infty$, and that $\int_0^s\frac{1-e^{-t}}t\dd t=\log s+\gamma+o(1)$. Conclude that $\int_0^\infty\rho(u)\dd u=e^\gamma=1.78107\ldots$
\item By Mertens's third theorem, $\sum\frac1n$ over the $y$-smooth $n$ equals $\prod_{p\le y}(1-\frac1p)^{-1}\sim e^\gamma\log y$. Assuming that $\Psi(t,y)\approx t\rho(\frac{\log t}{\log y})$,
show that this is exactly what (d) predicts.
\end{enumerate}
\end{exercise}

\begin{exercise}[$\star\star\star$]\label{pb:buchstab}
\emph{Rough numbers and Buchstab's function.}
\begin{enumerate}[label=(\alph*)]
\item Show that $\Phi(x,y)=1+\pi(x)-\pi(y)$ for $\sqrt x\le y\le x$, and that for $x^{1/3}\le y<\sqrt x$,
\[
\Phi(x,y)=1+\pi(x)-\pi(y)+\#\{(p,q):y<p\le q,\ pq\le x\}.
\]
\item Let $1<u\le3$. Show that $\frac{\log y}x\Phi(x,y)\to\omega(u)$ when $y=x^{1/u}$, where $\omega(u)=\frac1u$ for $1<u\le2$ and $\omega(u)=\frac{1+\log(u-1)}u$ for $2\le u\le3$.
\item Check that $(u\omega(u))'=\omega(u-1)$ for $2<u<3$ (Buchstab's equation, which defines $\omega$ for all $u>1$). Show that $\omega(2)<e^{-\gamma}<\omega(3)$.
\item Independent sieving by the primes up to $y$ would predict $\Phi(x,y)\approx x\prod_{p\le y}(1-\frac1p)\sim e^{-\gamma}\frac x{\log y}$, that is $\omega\equiv e^{-\gamma}$. Where on $]1,3]$ is $\omega$ smallest, and by what factor
does the prediction overestimate there? Relate this to Legendre's failure (Volume~5, Exercise~8.2).
\end{enumerate}
\end{exercise}

\begin{figure}[ht]
\centering
\fig{ch13-buchstab}
\caption{Buchstab's function $\omega(u)$ (blue), computed from $\omega(u)=\frac1u$ on $[1,2]$ and Buchstab's equation, with the proportions $\frac{\log y}x\Phi(x,y)$ for
$x=10^8$ and $y=x^{1/u}$, $u=1.5,2,2.5,3,4$ (green), and the level $e^{-\gamma}$ predicted by independent sieving (red, dashed). The function crosses $e^{-\gamma}$
between $2$ and $3$, again near $3.2$ and $4$, and is within $10^{-4}$ of it from $u=4$ on; the data are off by a term of order $\frac1{\log y}$ (Problem~\ref{pb:buchstab}).}
\label{fig:buchstab}
\end{figure}

\begin{exercise}[$\star\star$]\label{pb:polymertens}
\emph{Mertens's theorem for polynomials.} Let $N_k$ be the number of monic irreducible polynomials of degree $k$ over $\mathbb F_q$, so that
$\sum_{k\mid m}kN_k=q^m$ (Volume~2, Chapter~9). Let $E_n=\prod_{k\le n}(1-q^{-k})^{N_k}$.
\begin{enumerate}[label=(\alph*)]
\item Show that $-\log E_n=H_n+\varepsilon_n$, where $H_n=\sum_{m\le n}\frac1m$ and $0\le\varepsilon_n\le4(1+\log n)q^{-n/2}$. Deduce that $E_n\sim\frac{e^{-\gamma}}n$.
\item A monic polynomial of degree $n$ with no irreducible factor of degree at most $\frac n2$ is irreducible, and the probability of that is $\frac1n+O(\frac{q^{-n/2}}n)$
(Problem~\ref{pb:randpoly}). Independent sieving would predict $E_{\lfloor n/2\rfloor}$. Show that the prediction exceeds the truth by a factor tending to $2e^{-\gamma}=1.1229\ldots$, the same factor as for the integers.
\item Show that $ne^\gamma E_n=1-\frac1{2n}+O(\frac1{n^2})$ for fixed $q$. Numerically, one finds $ne^\gamma E_n=0.952015$ for $q=5$, $n=10$. Compare.
\end{enumerate}
\end{exercise}

\begin{exercise}[$\star\star$]\label{pb:hawkins}
\emph{Hawkins's random primes.} In 1958 David Hawkins proposed a random sieve. Start with the integers $2,3,4,\dots$ and let $X_1=2$. Strike out each integer larger than $X_1$
independently with probability $\frac1{X_1}$; let $X_2$ be the least surviving integer larger than $X_1$; strike out each surviving integer larger than $X_2$ with
probability $\frac1{X_2}$, independently; and so on. The survivors $X_1<X_2<\dots$ are the \emph{Hawkins primes}. Let $Q_n=\prod_{X_k<n}(1-\frac1{X_k})$.
\begin{enumerate}[label=(\alph*)]
\item Show that, given everything that happened to the integers below $n$, the integer $n$ is a Hawkins prime with probability $Q_n$, and that
$\mathbb E[Q_{n+1}\mid\text{the past}]=Q_n(1-\frac{Q_n}n)$.
\item Let $q_2=1$ and $q_{n+1}=q_n-\frac{q_n^2}n$. Show that $\mathbb E[Q_n]\le q_n$ for all $n\ge2$.
\item Show that $\frac1{q_{n+1}}=\frac1{q_n}+\frac1{n-q_n}$, deduce that $\frac1{q_n}=\log n+O(1)$, and conclude that the expected number of Hawkins primes up to $x$ is at most
$(1+o(1))\frac x{\log x}$.
\item For the Hawkins primes, the probability of survival \emph{is} the sieve product $Q_n\approx\frac1{\log n}$. For the true primes, a sieve by the primes up to $\sqrt n$
decides primality, and $\prod_{p\le\sqrt n}(1-\frac1p)\sim\frac{2e^{-\gamma}}{\log n}$, which is not $\frac1{\log n}$. What makes the difference?
\end{enumerate}
\end{exercise}

\section{Laws, gaps and last digits}

The last three problems look at whole distributions rather than averages: the law of a repeated-factor count, the law of the gap in which a random integer falls,
and the law of the last digit of the next prime. Each time a simple model makes a prediction, and each time the data have something to say about it.

\begin{exercise}[$\star\star\star$]\label{pb:renyi}
\emph{The law of $\Omega(n)-\omega(n)$.} Volume~5, Exercise~9.1, shows that $\Omega(n)-\omega(n)$ is bounded on average. Here is its exact law (R\'enyi, 1955).
\begin{enumerate}[label=(\alph*)]
\item For complex $z$ with $|z|<2$, show that $f_z(n)=z^{\Omega(n)-\omega(n)}$ is multiplicative and $f_z(n)=\sum_{d\mid n}h_z(d)$, where $h_z(p)=0$ and $h_z(p^k)=z^{k-2}(z-1)$ for $k\ge2$.
Show that $\sum_d\frac{|h_z(d)|}d<\infty$.
\item Deduce from Wintner's theorem (Volume~5, Theorem~2.2) that
\[
\frac1x\sum_{n\le x}z^{\Omega(n)-\omega(n)}\to R(z)=\prod_p\Bigl(1-\frac1p\Bigr)\Bigl(1+\frac1{p-z}\Bigr).
\]
\item Show that for every $k\ge0$ the set $\{n:\Omega(n)-\omega(n)=k\}$ has a density $d_k$, and that $\sum_kd_kz^k=R(z)$. (Write $d_k$ for $n\le x$ as an integral over the
unit circle.) Check that $d_0=\frac6{\pi^2}$ and that the mean of the law is $\sum_p\frac1{p(p-1)}=0.77315\ldots$
\item Show that $2^kd_k\to\frac1{4C_2}=0.37869\ldots$, where $C_2=\prod_{p\ge3}\bigl(1-\frac1{(p-1)^2}\bigr)$ is the twin prime constant.
\item Numerically, one finds $d_1=0.200756$, $d_2=0.096302$, $d_3=0.047615$, and, among the integers up to $10^7$, the proportions $0.200758$, $0.096306$, $0.047616$.
\end{enumerate}
\end{exercise}

\begin{exercise}[$\star\star$]\label{pb:inspection}
\emph{The inspection paradox for prime gaps.} Let $g_n=p_{n+1}-p_n$.
\begin{enumerate}[label=(\alph*)]
\item Show that $\sum_{p_{n+1}\le x}g_n^2\ge(1+o(1))\,x\log x$.
\item Let $p_N$ be the largest prime up to $x$, choose an integer $m$ uniformly in $]2,p_N]$, and let $G$ be the length of the gap $]p_n,p_{n+1}]$ containing $m$. Show that
$\mathbb E[G]=\frac1{p_N-2}\sum_{n<N}g_n^2$, and deduce that $\mathbb E[G]\ge(1+o(1))\log x$, although the average gap is $\sim\log x$. Compare with the inspection paradox
of Volume~3, Proposition~9.10.
\item In Cram\'er's model (Volume~5, Chapter~10) the gaps near $x$ are roughly geometric with mean $\log x$. Show that for a geometric law $G_0$ with mean $\frac1q$,
$\frac{\mathbb E[G_0^2]}{\mathbb E[G_0]}=\frac{2-q}q$, so that the model predicts $\sum_{p_{n+1}\le x}g_n^2\sim2x\log x$.
\item Figure~\ref{fig:gaps} shows $\frac1{x\log x}\sum_{p_{n+1}\le x}g_n^2$ for $10^3\le x\le10^8$. What do you conclude? If one draws $m$ at random and takes the next prime,
which primes are favoured?
\end{enumerate}
\end{exercise}

\begin{figure}[ht]
\centering
\fig{ch13-gaps}
\caption{The second moment of the prime gaps, $\frac1{x\log x}\sum_{p_{n+1}\le x}(p_{n+1}-p_n)^2$, for $x$ from $10^3$ to $10^8$ (blue): $1.16$ at $10^3$, $1.52$ at $10^6$,
$1.63$ at $10^8$. Problem~\ref{pb:inspection} proves that the limit inferior is at least $1$; Cram\'er's model predicts $2$.}
\label{fig:gaps}
\end{figure}

\begin{exercise}[$\star\star$]\label{pb:digitpairs}
\emph{The last digit of the next prime.} Among the primes $7\le p<10^8$, a prime ending in $1$ is followed by a prime ending in $1$, $3$, $7$, $9$ with frequencies
$17.7\%$, $30.4\%$, $31.0\%$, $20.9\%$ (Figure~\ref{fig:digits}). Here is a model. Walk through the integers prime to $10$, whose last digits cycle as
$1\to3\to7\to9\to1$, and declare each of them ``prime'' independently with probability $q$.
\begin{enumerate}[label=(\alph*)]
\item Show that after a ``prime'' ending in $1$, the next ``prime'' ends in $3$, $7$, $9$, $1$ with probabilities $\pi_j=\frac{q(1-q)^{j-1}}{1-(1-q)^4}$, $j=1,2,3,4$.
\item Show that $\pi_j=\frac14+\frac{(5-2j)q}8+O(q^2)$. Which transition does the model make rarest?
\item Explain why $q$ should be about $\frac{2.5}{\log x}$ near $x$. With this $q$, averaged over the primes ending in $1$ below $10^8$, the model gives $31.1\%$, $26.6\%$, $22.8\%$, $19.5\%$
for the endings $3$, $7$, $9$, $1$. Which part of the bias does it explain?
\item The model badly underestimates the transition $1\to7$. The shortest gap that makes it is $6$. Using the Hardy--Littlewood prediction that pairs $p$, $p+h$ (both prime) have
density proportional to $\prod_{p\mid h,\,p>2}\frac{p-1}{p-2}$ (Volume~5, Chapter~10), explain what the model has forgotten.
\end{enumerate}
\end{exercise}

\begin{figure}[ht]
\centering
\fig{ch13-digits}
\caption{After a prime ending in $1$, the last digit of the next prime: the frequencies among the primes $7\le p<10^8$ (blue) and the predictions of the coin model of
Problem~\ref{pb:digitpairs} with $q=\frac{2.5}{\log p}$ (orange). Both make a repeated $1$ the rarest transition; only the primes favour $7$ as much as $3$.}
\label{fig:digits}
\end{figure}

\section{Hints}

\begin{hint}{pb:leastpf}
(a) Abel summation (Volume~5, Lemma~1.1): $\sum_{p\le x}p=x\pi(x)-\int_2^x\pi(t)\dd t$. (b) A composite $n$ has $P^-(n)\le\sqrt n$. (c) $p_n\sim n\log n$. (d) Substitute $s=t^2$ in
$\int\frac t{\log t}\dd t$, and use $\pi=\operatorname{Li}+E$ with Volume~5, Corollary~7.9.
\end{hint}

\begin{hint}{pb:largestpf}
(a) Write $n=pm$ with $p=P^+(n)$. (c) $\lfloor\frac xp\rfloor$ is the number of $k\ge1$ with $p\le\frac xk$. Then let $x\to\infty$ in each term $\frac{\log x}{x^2}\sum_{p\le x/k}p$ and dominate them by $\frac C{k^2}$.
(d) $A(n)-P^+(n)$ is a sum of at most $\log_2n$ primes, each at most the second largest prime factor of $n$.
\end{hint}

\begin{hint}{pb:radical}
(a) $g=f*\mu$ with $f(p^k)=p^{1-k}$; compute $g(p^k)$. (b) Every squarefull number is $a^2b^3$. (c) $\sum_{n\le x}\operatorname{rad}(n)=\sum_dg(d)\,d\sum_{m\le x/d}m$.
\end{hint}

\begin{hint}{pb:factorialdiv}
(c) Telescope: $\log(1+m)=\sum_{k=1}^m\log(1+\frac1k)$ with $m=\lfloor\frac np\rfloor$. Exchange the sums, and dominate as in Problem~\ref{pb:largestpf}.
\end{hint}

\begin{hint}{pb:centralbinom}
(a) $\lfloor2t\rfloor-2\lfloor t\rfloor\in\{0,1\}$. (c) For $k\le n^{1/4}$ the intervals are long; use Brun--Titchmarsh (Volume~5, Theorem~8.4) to dominate. The primes below
$n^{3/4}$ are negligible.
\end{hint}

\begin{hint}{pb:primezeta}
(a) Take logarithms in the Euler product (Volume~4, Theorem~10.3). (b) $\sum_{k\mid m}\mu(k)=[m=1]$. (c) $\log\zeta(\sigma)\le\zeta(\sigma)-1\le2^{1-\sigma}$ for $\sigma\ge3$.
\end{hint}

\begin{hint}{pb:superprimes}
(a) $p_m\sim m\log m$ with $m=p_n$. (b) Compare with $\sum\frac1{n\log^2n}$.
\end{hint}

\begin{hint}{pb:ramprimes}
(a) $f$ jumps up by $1$ at each prime and down by $1$ at each $2p$. (b) Compare $n=f(R_n)\sim\frac{R_n}{2\log R_n}$ with $2n=\pi(p_{2n})\sim\frac{p_{2n}}{\log p_{2n}}$.
\end{hint}

\begin{hint}{pb:pminusone}
(a) Split $q\le x^{1/4}$, where~\eqref{eq:bt13} applies to $d=q^2$, and $q>x^{1/4}$. (b) For fixed $y$, the condition ``$q^2\nmid p-1$ for all $q\le y$'' depends only on $p$
modulo $\prod_{q\le y}q^2$. (c) $\frac{\varphi(m)}m=\sum_{d\mid m}\frac{\mu(d)}d$; treat $d\le y$, $y<d\le\sqrt x$ and $d>\sqrt x$ separately. For the last range put $e=\frac{p-1}d$.
\end{hint}

\begin{hint}{pb:titchmarsh}
(b) $\tau(m)\le2\#\{d\mid m:d\le\sqrt m\}$. (c) $\frac m{\varphi(m)}=\sum_{d\mid m}\frac{\mu(d)^2}{\varphi(d)}$, and as in Problem~\ref{pb:pminusone}; you need $\varphi(d)\ge c_\varepsilon d^{1-\varepsilon}$,
which follows from $\frac d{\varphi(d)}\le\tau(d)$ and Volume~5, Exercise~2.6.
\end{hint}

\begin{hint}{pb:menon}
(a) $\mathbb F_p^\times$ is cyclic of order $p-1$. (b) Burnside's lemma: the orbit of $x$ is determined by $\gcd(x,n)$; the fixed points of $a$ are the $x$ with $(a-1)x\equiv0$.
(c) $\gcd(n,p-1)$ depends only on $p\bmod n$. (d) $\Phi_d$ has $\varphi(d)$ roots modulo $p\nmid n$ if $d\mid p-1$, and none otherwise.
\end{hint}

\begin{hint}{pb:omegapoly}
(b) Exchange: $\sum_n\omega(n^2+1)=\sum_p\#\{n\le x:p\mid n^2+1\}$. (c) $\sum_p\frac{\chi(p)}p$ converges (Volume~5, Corollary~4.5). (d) Work with $\omega_y(m)$, the number of prime
factors $p\le y=x^{1/4}$ of $m$, and compute $\sum_n\omega_y(n^2+1)^2$ by the Chinese remainder theorem.
\end{hint}

\begin{hint}{pb:rankin}
(a) For a $y$-smooth $n\le x$, $1\le(\frac xn)^\sigma$. (b) Split $\log\prod(1-p^{-\sigma})^{-1}=\sum_pp^{-\sigma}+\sum_p\sum_{k\ge2}\frac{p^{-k\sigma}}k$, and note $\sigma\ge\frac23$.
(d) Compare $\prod_{p\le y}(1-\frac1p)^{-1}\sim e^\gamma\log y$ with the true size.
\end{hint}

\begin{hint}{pb:dickmanint}
(a) Both sides have the same derivative and agree at $u=1$. If $u_0$ were the first zero of $\rho$, the identity at $u_0$ would fail. (b) Integrate $e^{-su}(u\rho(u))'$ by parts.
(d) Split $\int_0^s=\int_0^1+\int_1^s$ and use $\gamma=\int_0^1\frac{1-e^{-t}}t\dd t-\int_1^\infty\frac{e^{-t}}t\dd t$, which follows from $\Gamma'(1)=-\gamma$ (Volume~4, Exercise~5.7).
(e) $\sum\frac1n=\int_{1^-}^\infty\frac{\mathrm d\Psi(t,y)}t=\int_1^\infty\frac{\Psi(t,y)}{t^2}\dd t$.
\end{hint}

\begin{hint}{pb:buchstab}
(a) A $y$-rough $n\le x$ with $y\ge x^{1/3}$ has at most two prime factors. (b) Count the pairs as $\sum_{y<p\le\sqrt x}(\pi(\frac xp)-\pi(p)+1)$, and turn
$\sum_{x^a<p\le x^b}\frac1pF(\frac{\log p}{\log x})$ into $\int_a^bF(t)\frac{\mathrm dt}t$ with Mertens's second theorem.
\end{hint}

\begin{hint}{pb:polymertens}
(a) Expand $-\log(1-q^{-k})=\sum_j\frac{q^{-jk}}j$ and group the terms with $jk=m\le n$. The others have $j\ge2$ and $(j-1)k\ge\frac n2$. (c) $H_n=\log n+\gamma+\frac1{2n}+O(\frac1{n^2})$.
\end{hint}

\begin{hint}{pb:hawkins}
(a) $Q_{n+1}=Q_n(1-\frac1n)$ if $n$ is a Hawkins prime, $Q_{n+1}=Q_n$ otherwise. (b) Jensen, and $t\mapsto t-\frac{t^2}n$ increases on $[0,1]$. (c) $0<q_n\le1$.
\end{hint}

\begin{hint}{pb:renyi}
(b) Compute the Euler factor $1+\sum_{k\ge2}\frac{h_z(p^k)}{p^k}$. (c) Cauchy's formula on $|z|=1$ and dominated convergence. (d) The factor $p=2$ of $R$ has a pole at $z=2$; the rest
is analytic in $|z|<3$.
\end{hint}

\begin{hint}{pb:inspection}
(a) Cauchy--Schwarz, $\sum g_n\le x$ with $\pi(x)$ terms. (b) Each $m$ in a gap of length $g$ sees $G=g$, and there are $g$ such $m$. (c) $\mathbb E[G_0]=\frac1q$, $\mathbb E[G_0^2]=\frac{2-q}{q^2}$.
\end{hint}

\begin{hint}{pb:digitpairs}
(a) Sum over the number of complete turns of the cycle. (c) Among the integers near $x$, a proportion $\frac4{10}$ is prime to $10$. (d) Compare the weights of $h=2$, $4$ and $6$.
\end{hint}

\section{Solutions}

\begin{solution}{pb:leastpf}
\idea Weighting the primes by $p$ makes the largest primes count most, so the sum behaves like $x\cdot\pi(x)$ corrected by an integral. The least prime factor of a composite is tiny,
so the composites do not matter.

\writeup (a) By Abel summation (Volume~5, Lemma~1.1) with $a_p=p$, or directly, $\sum_{p\le x}p=x\pi(x)-\int_2^x\pi(t)\dd t$. By the prime number theorem $x\pi(x)\sim\frac{x^2}{\log x}$.
For the integral, fix $\varepsilon>0$; for $t\ge t_0$, $\pi(t)\le(1+\varepsilon)\frac t{\log t}$, and $\frac t{\log t}\le\frac t{\log\sqrt x}=\frac{2t}{\log x}$ for $t\ge\sqrt x$. So
\[
\int_2^x\pi(t)\dd t\le O(t_0^2)+\int_2^{\sqrt x}t\dd t+(1+\varepsilon)\int_{\sqrt x}^x\frac t{\log t}\dd t.
\]
On $[x^{1-\varepsilon},x]$, $\frac1{\log t}\le\frac1{(1-\varepsilon)\log x}$, and on $[\sqrt x,x^{1-\varepsilon}]$ the integrand is at most $\frac{2t}{\log x}$, with
$\int_{\sqrt x}^{x^{1-\varepsilon}}t\dd t\le\frac12x^{2-2\varepsilon}$. Hence $\int_2^x\pi(t)\dd t\le\frac{1+\varepsilon}{1-\varepsilon}\frac{x^2}{2\log x}+o(\frac{x^2}{\log x})$, and the same argument
with $(1-\varepsilon)$ and $\frac1{\log t}\ge\frac1{\log x}$ gives the matching lower bound. So $\int_2^x\pi(t)\dd t\sim\frac{x^2}{2\log x}$ and $\sum_{p\le x}p\sim\frac{x^2}{\log x}-\frac{x^2}{2\log x}=\frac{x^2}{2\log x}$.

(b) If $n$ is prime, $P^-(n)=n$; if $n$ is composite, $P^-(n)^2\le n$, so $P^-(n)\le\sqrt x$. Hence $0\le\sum_{2\le n\le x}P^-(n)-\sum_{p\le x}p\le x\cdot\sqrt x=o\bigl(\frac{x^2}{\log x}\bigr)$.

(c) The sum of the first $n$ primes is $\sum_{p\le p_n}p\sim\frac{p_n^2}{2\log p_n}$ by (a). Since $p_n\sim n\log n$ (Volume~5, Exercise~5.4), $\log p_n\sim\log n$ and the sum is
$\sim\frac{n^2\log^2n}{2\log n}=\frac12n^2\log n$.

(d) Heuristically $\mathrm d\pi(t)\approx\frac{\mathrm dt}{\log t}$, so $\sum_{p\le x}p\approx\int_2^x\frac t{\log t}\dd t$. With $s=t^2$, $\mathrm ds=2t\dd t$ and $\log t=\frac12\log s$, so
$\frac{t\dd t}{\log t}=\frac{\mathrm ds}{\log s}$ and the integral is $\operatorname{Li}(x^2)$ up to a constant. Precisely, write $\pi=\operatorname{Li}+E$ with $E(t)=O(te^{-c\sqrt{\log t}})$ (Volume~5,
Corollary~7.9). Then
\[
\sum_{p\le x}p=\int_{2^-}^xt\dd\pi(t)=\int_2^x\frac t{\log t}\dd t+\Bigl[tE(t)\Bigr]_2^x-\int_2^xE(t)\dd t.
\]
The first integral is $\operatorname{Li}(x^2)-\operatorname{Li}(4)$. The boundary term is $O(x^2e^{-c\sqrt{\log x}})$, and so is the last integral, since $te^{-c\sqrt{\log t}}$ is increasing for
large $t$ and $\int_2^x\le x\cdot\max$. So $\sum_{p\le x}p=\operatorname{Li}(x^2)+O(x^2e^{-c\sqrt{\log x}})$. The expansion $\operatorname{Li}(y)=\frac y{\log y}(1+\frac1{\log y}+\dots)$ at
$y=x^2$ gives $\frac{x^2}{2\log x}(1+\frac1{2\log x}+\dots)$; at $x=10^7$ the correction $\frac1{2\log x}$ is $3.1\%$, which explains the ratio $1.0326$.

\tools{Volume~5, Lemma~1.1 (Abel summation), Theorem~5.5 (prime number theorem), Corollary~7.9 (its error term), Exercise~5.4 ($p_n\sim n\log n$).}
\leads Weighting by $p$ moves the action to the top of the range, which is why the crude form $\frac{x^2}{2\log x}$ is so poor: it inherits the relative error $\frac1{\log x}$
of $\frac x{\log x}$. The same mechanism drives the next problem, where the weight is the largest prime factor.
\end{solution}

\begin{solution}{pb:largestpf}
\idea The largest prime factor is large exactly when $n$ is a prime times a small cofactor. Counting by the cofactor $k=\frac np$ turns the sum into $\sum_k\sum_{p\le x/k}p$,
and each inner sum is $\frac{x^2}{2k^2\log x}$ to first order: the $\frac{\pi^2}{12}$ is $\frac12\zeta(2)$.

\writeup (a) Every $n\ge2$ is uniquely $n=pm$ with $p=P^+(n)$ and $m\le\frac xp$ having no prime factor larger than $p$; conversely every such pair gives an $n\le x$ with $P^+(n)=p$.
So $\sum_{2\le n\le x}P^+(n)=\sum_{p\le x}p\,\Psi(\frac xp,p)$.

(b) For $p\le\sqrt x$, $p\,\Psi(\frac xp,p)\le p\cdot\frac xp=x$, and there are at most $\sqrt x$ such primes: total $O(x^{3/2})$. For $p>\sqrt x$, every $m\le\frac xp<\sqrt x<p$
has all its prime factors below $p$, so $\Psi(\frac xp,p)=\lfloor\frac xp\rfloor$.

(c) Let $T(y)=\sum_{p\le y}p$. Since $\lfloor\frac xp\rfloor=\#\{k\ge1:k\le\frac xp\}=\#\{k\ge1:p\le\frac xk\}$,
\[
\sum_{\sqrt x<p\le x}p\Bigl\lfloor\frac xp\Bigr\rfloor=\sum_{k\ge1}\ \sum_{\sqrt x<p\le x/k}p=\sum_{1\le k<\sqrt x}\bigl(T(\tfrac xk)-T(\sqrt x)\bigr),
\]
because the inner sum is empty when $\frac xk\le\sqrt x$. The terms $T(\sqrt x)\le x$ contribute at most $x^{3/2}$. For the rest, by Problem~\ref{pb:leastpf}(a), for each fixed $k$,
$\frac{\log x}{x^2}T(\frac xk)\to\frac1{2k^2}$. For a bound, $T(y)\le y\pi(y)\le C\frac{y^2}{\log y}$ (Chebyshev), and for $k<\sqrt x$, $\log\frac xk\ge\frac12\log x$, so
$\frac{\log x}{x^2}T(\frac xk)\le\frac{2C}{k^2}$. By dominated convergence for series,
\[
\frac{\log x}{x^2}\sum_{1\le k<\sqrt x}T\Bigl(\frac xk\Bigr)\to\sum_{k\ge1}\frac1{2k^2}=\frac{\pi^2}{12}.
\]
With (a) and (b), $\sum_{2\le n\le x}P^+(n)=\frac{\pi^2}{12}\frac{x^2}{\log x}(1+o(1))+O(x^{3/2})$, which is the claim.

(d) Write $n=P^+(n)\,m$. Then $A(n)-P^+(n)=A(m)$ is a sum of $\Omega(m)\le\log_2n$ primes, each at most $P^+(m)$. If $m>1$, $P^+(m)$ is a prime factor of $n$ other than the
largest one (counted with multiplicity), so $P^+(m)^2\le P^+(m)P^+(n)\le n$. Hence $0\le A(n)-P^+(n)\le\sqrt n\log_2n$, and summing,
$0\le\sum_{n\le x}A(n)-\sum_{n\le x}P^+(n)\le x^{3/2}\log_2x=o\bigl(\frac{x^2}{\log x}\bigr)$.

(e) In (c), $T(\frac xk)\approx\frac{x^2}{2k^2\log(x/k)}$, and $\log\frac xk<\log x$: each term is larger than its limit, by a factor $1+\frac{\log k}{\log x}+\dots$ The
corrections are of relative size $\frac1{\log x}$, about $6\%$ at $x=10^7$, so the convergence is logarithmically slow.

\tools{Volume~5, Theorem~5.5 and Proposition~1.2; Problem~\ref{pb:leastpf}; $\zeta(2)=\frac{\pi^2}6$ (Volume~4).}
\leads Alladi and Erd\H os (1977) found this asymptotic formula, in the equivalent form ``$A(n)$ has average order $\frac{\pi^2}6\frac n{\log n}$'', and extended it to the $k$-th largest prime
factor. The average of $\frac{\log P^+(n)}{\log n}$, a different question, leads to the Golomb--Dickman constant and to Dickman's function of Section~3.
\end{solution}

\begin{solution}{pb:radical}
\idea The radical forgets repeated prime factors, and only numbers with repeated factors make $\frac{\operatorname{rad}(n)}n$ differ from $1$. Squarefull numbers are rare,
about $\sqrt x$ of them, so the convolution trick of Volume~5, Chapter~2, works with a power saving.

\writeup (a) $f(n)=\frac{\operatorname{rad}(n)}n$ is multiplicative with $f(p^k)=p^{1-k}$. Möbius inversion gives a multiplicative $g=f*\mu$ with $g(p^k)=f(p^k)-f(p^{k-1})$:
\[
g(p)=1-1=0,\qq g(p^k)=p^{1-k}-p^{2-k}=-(p-1)p^{1-k}\quad(k\ge2).
\]
So $g(d)=0$ unless $d$ is squarefull, and then $|g(d)|=\prod_{p^k\parallel d}(p-1)p^{1-k}\le1$, since $p-1\le p^{k-1}$ for $k\ge2$.

(b) Every squarefull $d$ can be written $d=a^2b^3$ (put one or three copies of each prime with odd exponent into $b^3$). For fixed $b$ there are at most $\sqrt{y/b^3}$ values of
$a$, so the count is at most $\sqrt y\sum_bb^{-3/2}=\zeta(\frac32)\sqrt y$.

(c) Since $\operatorname{rad}(n)=n\sum_{d\mid n}g(d)$, writing $n=dm$,
\[
\begin{aligned}\sum_{n\le x}\operatorname{rad}(n)&=\sum_{d\le x}g(d)\,d\sum_{m\le x/d}m=\sum_{d\le x}g(d)\,d\Bigl(\frac{x^2}{2d^2}+O\Bigl(\frac xd\Bigr)\Bigr)\\&=\frac{x^2}2\sum_{d\le x}\frac{g(d)}d+O\Bigl(x\sum_{d\le x}|g(d)|\Bigr).\end{aligned}
\]
By (a) and (b), $\sum_{d\le x}|g(d)|\le\zeta(\frac32)\sqrt x$, and by partial summation
\[
\sum_{d>x}\frac{|g(d)|}d\le\sum_{\substack{d>x\\ d\text{ squarefull}}}\frac1d=O(x^{-1/2}).
\]
So
$\sum_{n\le x}\operatorname{rad}(n)=\frac{x^2}2G+O(x^{3/2})$ with $G=\sum_d\frac{g(d)}d$. The series converges absolutely, so it is an Euler product:
\[
G=\prod_p\Bigl(1-(p-1)\sum_{k\ge2}\frac{p^{1-k}}{p^k}\Bigr)=\prod_p\Bigl(1-\frac{(p-1)p^{-3}}{1-p^{-2}}\Bigr)=\prod_p\Bigl(1-\frac1{p(p+1)}\Bigr)=C.
\]
The mean value of $f=\mathbf1*g$ is $\sum_d\frac{g(d)}d=C$ by Wintner's theorem (Volume~5, Theorem~2.2). Numerically, one finds $\frac2{x^2}\sum_{n\le10^7}\operatorname{rad}(n)=0.704446$, against $C=0.704442$.

\tools{Volume~5, Theorem~2.2 (Wintner) and the convolution method of Chapter~2; Euler products (Volume~4, Theorem~10.3).}
\leads Both the mean of $\frac{\operatorname{rad}(n)}n$ and the mean of $\operatorname{rad}(n)$ relative to $\frac x2$ equal $C$: the weight $n$ only rescales by $\frac12$. So a typical radical
is a fixed fraction of $n$, and numbers with a small radical, the squarefull ones, are rare. The $abc$ conjecture asserts that when $a+b=c$ with coprime $a,b$, the radical
of $abc$ can be much smaller than $c$ only in a weak sense: $c\le K_\varepsilon\operatorname{rad}(abc)^{1+\varepsilon}$.
\end{solution}

\begin{solution}{pb:factorialdiv}
\idea A prime $p>\sqrt n$ divides $n!$ exactly $\lfloor\frac np\rfloor$ times, and there are about $\frac n{k\log n}-\frac n{(k+1)\log n}$ primes with $\lfloor\frac np\rfloor=k$. Each contributes
$\log(1+k)$; summing by parts gives $\sum_k\frac1k\log(1+\frac1k)$.

\writeup (a) By Legendre's formula $n!=\prod_{p\le n}p^{v_p(n!)}$ with $v_p(n!)=\sum_j\lfloor\frac n{p^j}\rfloor$, and a number $\prod p^{a_p}$ has $\prod(1+a_p)$ divisors.

(b) For $p\le\sqrt n$, $\log(1+v_p(n!))\le\log(1+n)$, and there are at most $\sqrt n$ such primes: $O(\sqrt n\log n)$, which is $o(\frac n{\log n})$. (Even $O(\sqrt n)$ with
Chebyshev's bound, but we do not need it.) For $p>\sqrt n$, $p^2>n$, so $v_p(n!)=\lfloor\frac np\rfloor$.

(c) Telescoping, $\log(1+m)=\sum_{k=1}^m\log(1+\frac1k)$. So
\[
\sum_{\sqrt n<p\le n}\log\Bigl(1+\Bigl\lfloor\frac np\Bigr\rfloor\Bigr)=\sum_{k\ge1}\log\Bigl(1+\frac1k\Bigr)\#\Bigl\{\sqrt n<p\le\frac nk\Bigr\},
\]
which equals
\[
\sum_{1\le k<\sqrt n}\log\Bigl(1+\frac1k\Bigr)\Bigl(\pi\Bigl(\frac nk\Bigr)-\pi(\sqrt n)\Bigr).
\]
The terms with $\pi(\sqrt n)$ add up to at most $\sqrt n\sum_{k<\sqrt n}\frac1k=O(\sqrt n\log n)$. For each fixed $k$, $\frac{\log n}n\pi(\frac nk)\to\frac1k$, and for $k<\sqrt n$,
$\frac{\log n}n\pi(\frac nk)\le\frac{\log n}n\cdot\frac{Cn}{k\log(n/k)}\le\frac{2C}k$. Since $\log(1+\frac1k)\le\frac1k$, the terms are dominated by $\frac{2C}{k^2}$, and
\[
\frac{\log n}n\log\tau(n!)\to\sum_{k\ge1}\frac1k\log\Bigl(1+\frac1k\Bigr)=C.
\]
Finally $e^C=\prod_k(1+\frac1k)^{1/k}=2\cdot(\frac32)^{1/2}(\frac43)^{1/3}\cdots$

(d) The convergence is slow for the same reason as in Problem~\ref{pb:largestpf}: $\pi(\frac nk)\approx\frac n{k\log(n/k)}(1+\frac1{\log(n/k)})$, so the true sum is larger than its limit
by terms of relative order $\frac1{\log n}$; the excess over $C$ is $0.33$ at $n=10^4$ and $0.16$ at $n=10^7$, about $\frac3{\log n}$ in both cases.

\tools{Legendre's formula (Volume~1); Volume~5, Theorem~5.5 and Proposition~1.2; dominated convergence for series (Volume~0).}
\leads Erd\H os, Graham, Ivi\'c and Pomerance (1996) gave a full asymptotic expansion of $\log\tau(n!)$ in powers of $\frac1{\log n}$, and studied how fast $\tau((n+k)!)$ outgrows $\tau(n!)$.
\end{solution}

\begin{solution}{pb:centralbinom}
\idea For a large prime $p$, whether $p$ divides $\binom{2n}n$ depends only on the fractional part of $\frac np$: it does when there is a carry in $n+n$ in base $p$. The primes for
which this happens fill a union of intervals, whose total "prime mass" is $\sum_k\bigl(\frac1{k+1/2}-\frac1{k+1}\bigr)=2\log2$.

\writeup (a) By Legendre, $v_p(\binom{2n}n)=\sum_j\bigl(\lfloor\frac{2n}{p^j}\rfloor-2\lfloor\frac n{p^j}\rfloor\bigr)$, and for $p^2>2n$ only $j=1$ survives. With $t=\frac np=\lfloor t\rfloor+\{t\}$,
$\lfloor2t\rfloor-2\lfloor t\rfloor=\lfloor2\{t\}\rfloor$, which is $1$ if $\{t\}\ge\frac12$ and $0$ otherwise.

(b) $\{\frac np\}\ge\frac12$ means $k+\frac12\le\frac np<k+1$ with $k=\lfloor\frac np\rfloor\ge0$, that is $\frac n{k+1}<p\le\frac n{k+1/2}$.

(c) Let $I_k=\bigl]\frac n{k+1},\frac n{k+1/2}\bigr]$, of length $\ell_k=n\bigl(\frac1{k+1/2}-\frac1{k+1}\bigr)=\frac n{(2k+1)(k+1)}$. The primes $\le\sqrt{2n}$ contribute at most $\sqrt{2n}$.
For each fixed $k$, the prime number theorem gives $\#(I_k\cap\mathcal P)=\pi(\frac n{k+1/2})-\pi(\frac n{k+1})\sim\frac{\ell_k}{\log n}$. For $k\le n^{1/4}$, $\ell_k\ge\frac n{2(k+1)^2}\ge\frac{\sqrt n}3$
for large $n$, so $\log\ell_k\ge0.4\log n$ and Brun--Titchmarsh (Volume~5, Theorem~8.4) gives $\#(I_k\cap\mathcal P)\le\frac{3\ell_k}{\log\ell_k}\le\frac{8\ell_k}{\log n}$. As $\frac{\ell_k}n\le\frac1{(k+1)^2}$
is summable, dominated convergence gives
\[
\frac{\log n}n\sum_{k\le n^{1/4}}\#(I_k\cap\mathcal P)\to\sum_{k\ge0}\frac1{(2k+1)(k+1)}=2\sum_{k\ge0}\Bigl(\frac1{2k+1}-\frac1{2k+2}\Bigr)=2\log2.
\]
The intervals with $k>n^{1/4}$ lie below $n^{3/4}$ and contain at most $\pi(n^{3/4})=o(\frac n{\log n})$ primes. Hence $\omega\bigl(\binom{2n}n\bigr)\sim2\log2\cdot\frac n{\log n}=\log4\cdot\frac n{\log n}$.

(d) For $n=5\cdot10^6$, $\log4\cdot\frac n{\log n}=449\,370$, and $453\,587$ is $1.0094$ times this. (The ratio is above $1$ for the reason met in the previous problems: the
intervals $I_k$ sit at $\frac nk$, where $\log$ is smaller than $\log n$.)

\tools{Legendre's formula and Kummer's carries (Volume~1); Volume~5, Theorems~5.5 and~8.4.}
\leads The primes of $]n,2n]$ ($k=0$) give $\frac n{\log n}$, and those of $]\frac n2,\frac{2n}3]$ ($k=1$) give a sixth of that. Kummer's theorem says more: $v_p(\binom{2n}n)$ is the number of carries in $n+n$ in base $p$,
which for small $p$ can be large; but small primes are few.
\end{solution}

\begin{solution}{pb:primezeta}
\idea The logarithm of the Euler product is a sum of prime zeta values at multiples of $s$; Möbius inversion over the multiples solves for $P(s)$. The values $\zeta(ks)$ are
extremely close to $1$, so the series converges like $4^{-k}$ for $s=2$.

\writeup (a) For $s>1$, $\zeta(s)=\prod_p(1-p^{-s})^{-1}$ (Volume~4, Theorem~10.3), so $\log\zeta(s)=\sum_p\sum_{j\ge1}\frac{p^{-js}}j=\sum_{j\ge1}\frac{P(js)}j$, all terms being positive.

(b) Using (a) at $ks$,
\[
\sum_{k\ge1}\frac{\mu(k)}k\log\zeta(ks)=\sum_{k\ge1}\sum_{j\ge1}\frac{\mu(k)}{kj}P(jks)=\sum_{m\ge1}\frac{P(ms)}m\sum_{k\mid m}\mu(k)=P(s).
\]
The rearrangement is legitimate because $\sum_{k,j}\frac{P(jks)}{kj}=\sum_m\frac{\tau(m)P(ms)}m$, and $P(ms)\le\zeta(ms)-1\le2^{1-ms}$ for $ms\ge3$, so the series converges.

(c) For $\sigma\ge3$, $\log\zeta(\sigma)\le\zeta(\sigma)-1\le2^{-\sigma}+\int_2^\infty t^{-\sigma}\dd t=2^{-\sigma}\bigl(1+\frac2{\sigma-1}\bigr)\le2^{1-\sigma}$. For $k\ge K+1$ and $s\ge2$,
$ks\ge4$, so the error is at most
\[
\sum_{k>K}\frac{2^{1-ks}}k\le\frac{2^{1-(K+1)s}}{(K+1)(1-2^{-s})}\le\frac83\cdot\frac{2^{-(K+1)s}}{K+1}\le\frac4K2^{-(K+1)s}.
\]
For $s=2$ and $K=20$ this is $0.2\cdot2^{-42}<5\cdot10^{-14}$.

(d) By the prime number theorem, $\sum_{p>x}p^{-2}=\int_x^\infty t^{-2}\dd\pi(t)\sim\int_x^\infty\frac{\mathrm dt}{t^2\log t}\sim\frac1{x\log x}$, which is $6.2\cdot10^{-9}$ at $x=10^7$. And indeed
$0.452247420041-0.452247414181=5.86\cdot10^{-9}$. (The slightly smaller true value comes from the next term of $\frac1{x\log x}(1-\frac1{\log x}+\dots)$.)

\tools{Volume~4, Theorem~10.3 (Euler product) and Corollary~10.4 (Möbius inversion); Volume~5, Theorem~5.5.}
\leads The same inversion computes any sum $\sum_pf(p)$ for which $\sum_nf(n)$ has an Euler product: Mertens's constant and the twin prime constant are computed this way to hundreds of digits.
Summing the primes directly is hopeless beyond a few digits: the tail decays only like $\frac1{x\log x}$.
\end{solution}

\begin{solution}{pb:superprimes}
\idea Applying the prime number theorem twice: $\pi(\pi(x))$ loses one factor $\log x$ at each step.

\writeup (a) With $m=p_n\sim n\log n$, $p_m\sim m\log m$ and $\log m\sim\log n$, so $p_{p_n}\sim n\log^2n$. Next, $\pi(x)\to\infty$ and $\pi(y)\sim\frac y{\log y}$, so
$\pi(\pi(x))\sim\frac{\pi(x)}{\log\pi(x)}$; since $\log\pi(x)=\log x-\log\log x+o(1)\sim\log x$, $\pi(\pi(x))\sim\frac x{\log^2x}$.

(b) By (a), $\frac1{p_{p_n}}\le\frac2{n\log^2n}$ for large $n$, and $\sum\frac1{n\log^2n}$ converges (compare with $\int\frac{\mathrm dt}{t\log^2t}=-\frac1{\log t}$). The partial sum up to $10^8$ stops at
$n=397\,557$, and the tail is about $\sum_{n>N}\frac1{n\log^2n}\approx\frac1{\log N}=0.078$. (This overestimates a little, since $p_{p_n}$ is larger than $n\log^2n$; the sum
of the series is a little above $1$.)

(c) $\pi(10^8)=5\,761\,455$, and $\frac y{\log y}(1+\frac1{\log y})$ at $y=5\,761\,455$ gives $393\,900$, within $1\%$ of $397\,557$, against $35\%$ for $\frac x{\log^2x}=294\,700$.

\tools{Volume~5, Theorem~5.5 and Exercise~5.4.}
\leads Dressler and Parker (1975) showed with a computer-aided argument that every integer larger than $96$ is a sum of distinct primes of the form $p_{p_n}$: a thin set of primes,
of density $\frac1{\log^2x}$, is still a basis in this weak sense.
\end{solution}

\begin{solution}{pb:ramprimes}
\idea The function $f(x)=\pi(x)-\pi(\frac x2)$ is a random-looking staircase that drifts upwards like $\frac x{2\log x}$. The last time it is below $n$ is followed by an upward
step, and upward steps happen only at primes.

\writeup (a) $f$ is a right-continuous step function, increasing by $1$ at each prime $x=p$ and decreasing by $1$ at each $x=2p$; an up-step and a down-step never occur at the same
point, since a prime is never twice a prime. By the prime number theorem $f(x)\sim\frac x{2\log x}\to\infty$, so $S=\{x>0:f(x)<n\}$ is bounded;
it contains $]0,2[$, where $f=0$. Let $R_n=\sup S$. For $x>R_n$, $f(x)\ge n$, and by right continuity $f(R_n)\ge n$. There are points of $S$ arbitrarily close to $R_n$ from the
left, so $f(R_n^-)<n\le f(R_n)$: $f$ jumps up at $R_n$, so $R_n$ is a prime, the jump is $1$, and $f(R_n)=n$, $f(R_n^-)=n-1$. Clearly $R_n$ is the least real with the required property.

(b) $R_n\to\infty$ since $f$ is bounded on bounded sets. From $f(R_n)=n$, $n\sim\frac{R_n}{2\log R_n}$. Also $2n=\pi(p_{2n})\sim\frac{p_{2n}}{\log p_{2n}}$. Hence
$\frac{R_n}{\log R_n}\sim\frac{p_{2n}}{\log p_{2n}}$. Taking logarithms, $\log R_n-\log\log R_n=\log p_{2n}-\log\log p_{2n}+o(1)$, so $\log R_n\sim\log p_{2n}$, and then
$R_n\sim p_{2n}$.

(c) With $L=\log x$, $\pi(x)\approx\frac xL(1+\frac1L)$ and $\pi(\frac x2)\approx\frac x{2(L-\log2)}(1+\frac1L)\approx\frac x{2L}(1+\frac{1+\log2}L)$, so
$f(x)\approx\frac x{2L}(1+\frac{1-\log2}L)$. Equating $f(R_n)=n$ with $\pi(p_{2n})=2n$, that is $\frac{p_{2n}}{2L}(1+\frac1L)=n$, and ignoring the difference between $\log R_n$ and $\log p_{2n}$
in second-order terms,
\[
\frac{R_n}{p_{2n}}\approx\frac{\log R_n}{\log p_{2n}}\Bigl(1+\frac{\log2}{\log p_{2n}}\Bigr)\approx1+\frac{\log2}{\log p_{2n}}+\frac{\log(R_n/p_{2n})}{\log p_{2n}}.
\]
At $n=10^4$, $\log p_{2n}=12.32$ and $1+\frac{\log2}{12.32}=1.056$; including the last term, $1.056+\frac{0.074}{12.32}=1.062$. The true ratio is $\frac{242\,057}{224\,737}=1.077$. The
approximation explains the bulk of the excess, and also why it decays only like $\frac1{\log n}$.

\tools{Volume~5, Theorem~5.5; Bertrand's postulate (Volume~1).}
\leads Sondow (2009) also proved $p_{2n}<R_n$ for $n>1$ and conjectured $R_n<p_{3n}$, proved by Laishram (2010). Ramanujan's original paper (1919) found $R_1,\dots,R_5$ as the thresholds $2,11,17,29,41$.
\end{solution}

\begin{solution}{pb:pminusone}
\idea For a fixed set of small primes $q$, the condition on $p-1$ is a congruence condition on $p$, and the prime number theorem in progressions gives its density exactly. The
large $q$ are a nuisance, and Brun--Titchmarsh shows that they affect only a small proportion of primes.

\writeup (a) For $q\le x^{1/4}$, $d=q^2\le\sqrt x$, and~\eqref{eq:bt13} gives $\pi(x;q^2,1)\le\frac{5x}{q(q-1)\log x}$. Since $\sum_{q>y}\frac1{q(q-1)}\le\sum_{m>y}\frac1{m(m-1)}\le\frac2y$, these $q$
account for at most $\frac{10x}{y\log x}$ primes. For $q>x^{1/4}$, the number of $p\le x$ with $q^2\mid p-1$ is at most $\frac x{q^2}$, and
$\sum_{q>x^{1/4}}\frac x{q^2}\le x\sum_{m>x^{1/4}}\frac1{m(m-1)}\le2x^{3/4}$.

(b) Fix $y$ and let $M=\prod_{q\le y}q^2$. For $p>y$, whether $q^2\mid p-1$ for some $q\le y$ depends only on $p\bmod M$. By the Chinese remainder theorem, the classes
$a\in(\Z/M\Z)^\times$ with $q^2\nmid a-1$ for all $q\le y$ number $\prod_{q\le y}(\varphi(q^2)-1)$. By the prime number theorem in progressions (Volume~5, Theorem~5.8), each class
receives $\sim\frac{\pi(x)}{\varphi(M)}$ primes, so the primes $p\le x$ with $q^2\nmid p-1$ for all $q\le y$ number $\sim A_y\pi(x)$, where
\[
A_y=\prod_{q\le y}\frac{\varphi(q^2)-1}{\varphi(q^2)}=\prod_{q\le y}\Bigl(1-\frac1{q(q-1)}\Bigr).
\]
The set of $p$ with $p-1$ squarefree is contained in this one, and differs from it by at most the primes counted in (a). As $\pi(x)\sim\frac x{\log x}$, the proportion of
primes $p\le x$ with $p-1$ squarefree has $\limsup\le A_y$ and $\liminf\ge A_y-\frac{10}y$. Let $y\to\infty$: $A_y\to A$, and the proportion tends to $A$.

(c) Using $\frac{\varphi(m)}m=\sum_{d\mid m}\frac{\mu(d)}d$, split according to $d\le y$, $y<d\le\sqrt x$, $d>\sqrt x$. First,
\[
\sum_{p\le x}\sum_{\substack{d\mid p-1\\d\le y}}\frac{\mu(d)}d=\sum_{d\le y}\frac{\mu(d)}d\pi(x;d,1)=\pi(x)\Bigl(\sum_{d\le y}\frac{\mu(d)}{d\varphi(d)}+o(1)\Bigr)
\]
by Volume~5, Theorem~5.8, $y$ being fixed. Second, by~\eqref{eq:bt13},
$\sum_{y<d\le\sqrt x}\frac1d\pi(x;d,1)\le\frac{5x}{\log x}\sum_{d>y}\frac1{d\varphi(d)}$, and the series converges since $\varphi(d)\ge\sqrt{d/2}$ (for odd $p$, $\varphi(p^k)=p^{k-1}(p-1)\ge p^{k/2}$, and
$\varphi(2^k)\ge2^{k/2}/\sqrt2$). Third, if $d\mid p-1$ with $d>\sqrt x$, then $e=\frac{p-1}d<\sqrt x$ and $\frac1d=\frac e{p-1}$; writing $p-1=em$,
\[
\sum_{p\le x}\sum_{\substack{d\mid p-1\\d>\sqrt x}}\frac1d\le\sum_{e<\sqrt x}e\sum_{m\le x/e}\frac1{em}\le\sqrt x(1+\log x)=o(\pi(x)).
\]
Hence $\frac1{\pi(x)}\sum_{p\le x}\frac{\varphi(p-1)}{p-1}$ lies within $6\sum_{d>y}\frac1{d\varphi(d)}+o(1)$ of $\sum_{d\le y}\frac{\mu(d)}{d\varphi(d)}$. Letting $y\to\infty$, the
limit is $\sum_d\frac{\mu(d)}{d\varphi(d)}=\prod_q\bigl(1-\frac1{q\varphi(q)}\bigr)=A$.

(d) $\mathbb F_p^\times$ is cyclic of order $p-1$, so it has $\varphi(p-1)$ generators: the proportion of primitive roots among the non-zero residues is $\frac{\varphi(p-1)}{p-1}$. Averaging over $p\le x$,
(c) says that a uniformly random prime $p\le x$ and a uniformly random non-zero residue modulo $p$ give a primitive root with probability tending to $A$. Artin's conjecture asks the same for a
\emph{fixed} residue such as $2$, and the data agree to four decimals.

\tools{Volume~5, Theorem~5.8 and Exercise~8.4 (Brun--Titchmarsh); Euler products (Volume~4); cyclicity of $\mathbb F_p^\times$ (Volume~1).}
\leads The coincidence of the two limits is no accident: $p-1$ squarefree and ``a random residue is a primitive root'' are both conditions prime by prime, $q^2\nmid p-1$ with probability $1-\frac1{q(q-1)}$
and ``not a $q$-th power when $q\mid p-1$'' with probability $1-\frac1{q(q-1)}$. Stephens (1969) proved (c) with an error term, and Hooley (1967) proved Artin's conjecture assuming the generalised
Riemann hypothesis.
\end{solution}

\begin{solution}{pb:titchmarsh}
\idea Pairing each divisor $d$ of $m$ with $\frac md$ shows that $\tau(m)$ is twice the number of small divisors; counting the primes $p\equiv1\pmod d$ for $d\le\sqrt x$ with Brun--Titchmarsh
then costs only a constant factor.

\writeup (a) $\frac d{\varphi(d)}=\prod_{p\mid d}\frac p{p-1}=\prod_{p\mid d}\bigl(1+\frac1{p-1}\bigr)=\sum_{e\mid d}\frac{\mu(e)^2}{\varphi(e)}$ (Volume~5, Exercise~2.2). Hence, writing $d=em$,
\[
\sum_{d\le y}\frac1{\varphi(d)}=\sum_{e\le y}\frac{\mu(e)^2}{e\varphi(e)}\sum_{m\le y/e}\frac1m\le(1+\log y)\sum_e\frac1{e\varphi(e)}=c(1+\log y),
\]
the last series converging as in Problem~\ref{pb:pminusone}(c).

(b) The divisors of $m$ come in pairs $(d,\frac md)$ with $d\le\sqrt m$ or $\frac md\le\sqrt m$, so $\tau(m)\le2\#\{d\mid m:d\le\sqrt m\}$. With $m=p-1<x$,
\[
\sum_{p\le x}\tau(p-1)\le2\sum_{d\le\sqrt x}\pi(x;d,1)\le\frac{10x}{\log x}\sum_{d\le\sqrt x}\frac1{\varphi(d)}\le\frac{10x}{\log x}c\Bigl(1+\frac12\log x\Bigr)=O(x)
\]
by~\eqref{eq:bt13} and (a).

(c) Now $\frac m{\varphi(m)}=\sum_{d\mid m}\frac{\mu(d)^2}{\varphi(d)}$. The argument of Problem~\ref{pb:pminusone}(c) applies. For $d\le y$, the main term is $\pi(x)\sum_{d\le y}\frac{\mu(d)^2}{\varphi(d)^2}$.
For $y<d\le\sqrt x$, the contribution is at most $\frac{5x}{\log x}\sum_{d>y}\frac1{\varphi(d)^2}$; the series converges because $\frac d{\varphi(d)}\le2^{\omega(d)}\le\tau(d)\le C d^{1/4}$
(Volume~5, Exercise~2.6), so $\varphi(d)\ge\frac1Cd^{3/4}$. For $d>\sqrt x$, put $e=\frac{p-1}d<\sqrt x$; since $\varphi(de)\le e\varphi(d)$, $\frac1{\varphi(d)}\le\frac e{\varphi(p-1)}\le\frac{Ce}{(p-1)^{3/4}}$, and
\[
\sum_{p\le x}\sum_{\substack{d\mid p-1\\d>\sqrt x}}\frac1{\varphi(d)}\le C\sum_{e<\sqrt x}e\sum_{m\le x/e}\frac1{(em)^{3/4}}\le4C\sum_{e<\sqrt x}e^{1/4}\Bigl(\frac xe\Bigr)^{1/4}\le4Cx^{3/4}.
\]
Letting $x\to\infty$ and then $y\to\infty$, the mean tends to $\sum_d\frac{\mu(d)^2}{\varphi(d)^2}=\prod_q\bigl(1+\frac1{(q-1)^2}\bigr)=B$.

(d) Pretend that $\tau(p-1)=2\#\{d\mid p-1:d\le\sqrt x\}$ and that $\pi(x;d,1)\approx\frac{\pi(x)}{\varphi(d)}$ for all $d\le\sqrt x$. Then
\[
\sum_{p\le x}\tau(p-1)\approx2\pi(x)\sum_{d\le\sqrt x}\frac1{\varphi(d)}\approx2\frac x{\log x}\cdot Z\log\sqrt x=Zx.
\]
The data, $1.916$, $1.927$, $1.932$, creep up towards $Z=1.9436$.

\tools{Volume~5, Exercises~2.2, 2.6 and~8.4, Theorem~5.8.}
\leads Two different constants appear: on average over primes, $\frac{p-1}{\varphi(p-1)}$ is $B=2.83$, larger than the mean $Z=1.94$ of $\frac n{\varphi(n)}$ over all integers, because $p-1$ is
always even and more often divisible by small primes. The asymptotic formula $\sum_{p\le x}\tau(p-1)\sim Zx$ is the \emph{Titchmarsh divisor problem}: Titchmarsh (1930) proved it under the
generalised Riemann hypothesis, Linnik (1961) unconditionally with his dispersion method, and Halberstam (1967) gave a short proof from the Bombieri--Vinogradov theorem, the result that
makes the heuristic of (d) rigorous on average over $d\le\sqrt x/\log^Bx$.
\end{solution}

\begin{solution}{pb:menon}
\idea The number of roots is a gcd, and averaging a gcd over a group of units is a job for Burnside's lemma. The answer $\tau(n)$ is the number of irreducible factors of $x^n-1$ over $\Q$,
and each factor contributes one root on average.

\writeup (a) $\mathbb F_p^\times$ is cyclic of order $p-1$, and in a cyclic group of order $N$ the equation $x^n=1$ has exactly $\gcd(n,N)$ solutions. ($0$ is never a root.)

(b) Two elements $x,x'$ of $\Z/n\Z$ are in the same orbit if and only if $\gcd(x,n)=\gcd(x',n)$. Indeed multiplication by a unit preserves $\gcd(x,n)$; conversely if
$\gcd(x,n)=\gcd(x',n)=d$, then $x=dv$, $x'=dv'$ with $v,v'$ units modulo $\frac nd$, and since $(\Z/n\Z)^\times\to(\Z/\frac nd\Z)^\times$ is onto, some unit $a$ modulo $n$
reduces to $v'v^{-1}$ modulo $\frac nd$, and then $ax\equiv x'\pmod n$. So the orbits correspond to the divisors $d$ of $n$, and there are $\tau(n)$ of them. The fixed points of $a$
are the $x$ with $(a-1)x\equiv0\pmod n$, and there are $\gcd(a-1,n)$ of them. Burnside's lemma (Volume~2) gives
$\tau(n)=\frac1{\varphi(n)}\sum_{a\in(\Z/n\Z)^\times}\gcd(a-1,n)$.

(c) For $p\nmid n$, $\gcd(n,p-1)=\gcd(n,a-1)$ where $a\equiv p\pmod n$. Ignoring the finitely many $p\mid n$,
\[
\frac1{\pi(x)}\sum_{p\le x}\gcd(n,p-1)=\sum_{a\in(\Z/n\Z)^\times}\gcd(a-1,n)\frac{\pi(x;n,a)}{\pi(x)}+o(1),
\]
which tends to
\[
\frac1{\varphi(n)}\sum_a\gcd(a-1,n)=\tau(n)
\]
by Volume~5, Theorem~5.8, and (b).

(d) For $p\nmid n$, $x^n-1$ is separable modulo $p$ (its derivative $nx^{n-1}$ has no common root with it), and its roots are the elements of $\mathbb F_p^\times$ of order dividing $n$.
The roots of $\Phi_d$ are those of order exactly $d$: there are $\varphi(d)$ of them if $d\mid p-1$, none otherwise. By Volume~5, Theorem~5.8, a proportion $\frac1{\varphi(d)}$ of primes satisfy
$d\mid p-1$, so $\Phi_d$ has on average $\varphi(d)\cdot\frac1{\varphi(d)}=1$ root. Summing over $d\mid n$ gives $\tau(n)$ again, and the identity
$\gcd(n,p-1)=\sum_{d\mid\gcd(n,p-1)}\varphi(d)$ is the bridge between (c) and (d).

\tools{Burnside's lemma (Volume~2); cyclotomic polynomials and cyclicity of $\mathbb F_p^\times$ (Volumes~1 and~2); Volume~5, Theorem~5.8.}
\leads This is a special case of a theorem that goes back to Kronecker and Frobenius: for any $f\in\Z[x]$, the average number of roots of $f$ modulo $p$ is the number of irreducible factors of
$f$ over $\Q$. The general proof needs the Chebotarev density theorem (Stevenhagen and Lenstra tell the story). Menon's identity has since been generalised in dozens of ways.
\end{solution}

\begin{solution}{pb:omegapoly}
\idea Exchange the order of summation: $\sum_n\omega(n^2+1)$ counts pairs $(n,p)$ with $p\mid n^2+1$, and for each $p$ there are $\rho(p)$ residue classes of $n$. The primes with
$\rho(p)=2$ and $\rho(p)=0$ balance out to an average of one class per prime, as for $\omega(n)$ itself. For the variance, the classes modulo $p$ and $q$ are independent by the
Chinese remainder theorem, exactly as in Tur\'an's proof for $\omega(n)$.

\writeup (a) $n^2\equiv-1\pmod2$ means $n$ odd: one class. For odd $p$, $-1$ is a square modulo $p$ if and only if $p\equiv1\pmod4$ (Euler's criterion, Volume~1), and then it has
two square roots. So $\rho(p)=2$ or $0$, that is $1+\chi(p)$.

(b) If $n\le x$ and $n^2+1$ had two prime factors larger than $x$, then $n^2+1\ge(\lfloor x\rfloor+1)^2>\lfloor x\rfloor^2+1\ge n^2+1$. So
\[
\sum_{n\le x}\omega(n^2+1)=\sum_{p\le x}\#\{n\le x:p\mid n^2+1\}+O(x).
\]
For $p\le x$ the $n\le x$ with $p\mid n^2+1$ fill $\rho(p)$ classes modulo $p$, so there are $\rho(p)(\frac xp+O(1))$ of them. The errors add up to $O(\pi(x))$, and the claim follows.

(c) $\sum_{p\le x}\frac{\rho(p)}p=\sum_{p\le x}\frac1p+\sum_{p\le x}\frac{\chi(p)}p$ (the term $p=2$ fits since $\chi(2)=0$). The first sum is $\log\log x+O(1)$ (Volume~5, Theorem~1.5),
and the second converges, by Volume~5, Corollary~4.5 applied to both classes modulo~$4$. So $\sum_{n\le x}\omega(n^2+1)=x\log\log x+O(x)$.

(d) Let $y=x^{1/4}$ and $\omega_y(m)$ the number of prime factors $p\le y$ of $m$. For $x\ge3$, $n^2+1\le x^2+1<y^9$, so $0\le\omega(n^2+1)-\omega_y(n^2+1)\le8$. Let
$S=\sum_{p\le y}\frac{\rho(p)}p=\log\log x+O(1)$. Then
\[
\sum_{n\le x}\omega_y(n^2+1)^2=\sum_{p,q\le y}\#\{n\le x:p\mid n^2+1,\ q\mid n^2+1\}.
\]
For $p\ne q$, the condition is $pq\mid n^2+1$, which by the Chinese remainder theorem defines $\rho(p)\rho(q)$ classes modulo $pq$, so the count is $\rho(p)\rho(q)(\frac x{pq}+O(1))$;
for $p=q$ it is $\rho(p)(\frac xp+O(1))$. Hence $\sum_n\omega_y^2\le xS^2+xS+O(\pi(y)^2)$, and similarly $\sum_n\omega_y=xS+O(\pi(y))$. Expanding the square,
\[
\sum_{n\le x}\bigl(\omega_y(n^2+1)-S\bigr)^2\le xS^2+xS-2S(xS-O(\pi(y)))+xS^2+O(\pi(y)^2)=xS+O(\sqrt x).
\]
Finally $\omega(n^2+1)-\log\log x=(\omega_y-S)+O(1)$, and $(a+b)^2\le2a^2+2b^2$ gives $\sum_{n\le x}(\omega(n^2+1)-\log\log x)^2=O(x\log\log x)$. By Chebyshev's inequality, for every
$\varepsilon>0$ the $n\le x$ with $|\omega(n^2+1)-\log\log x|>\varepsilon\log\log x$ number $O\bigl(\frac x{\varepsilon^2\log\log x}\bigr)=o(x)$; and for $\sqrt x<n\le x$,
$|\log\log n-\log\log x|\le\log2$. So $\omega(n^2+1)$ has normal order $\log\log n$.

(e) Yes: (c) says that the difference is bounded, and it is, at about $0.6$. Its limit is a constant built from the Meissel--Mertens constant, $\sum_p\frac{\chi(p)}p$, and the
density of the $n$ for which $n^2+1$ has a prime factor larger than $n$.

\tools{Euler's criterion (Volume~1); Volume~5, Theorem~1.5, Corollary~4.5, and the method of Tur\'an's variance (Chapter~9); Chebyshev's inequality (Volume~3).}
\leads The same argument works for any irreducible polynomial $f$, with $\rho_f(p)$ averaging to $1$ by the theorem of Kronecker and Frobenius met in Problem~\ref{pb:menon}; and the
Erd\H os--Kac theorem holds for $\omega(f(n))$ too. The variance here is only bounded above; one finds numerically it smaller than $\log\log x$, partly because the large prime factors
have been thrown away.
\end{solution}

\begin{solution}{pb:rankin}
\idea Counting the smooth numbers up to $x$ is hard; summing $(\frac xn)^\sigma$ over \emph{all} smooth numbers is easy, because it factors as an Euler product. Choosing $\sigma$ slightly less
than $1$ balances the gain $x^{\sigma-1}$ against the growth of the product.

\writeup (a) For a $y$-smooth $n\le x$, $(\frac xn)^\sigma\ge1$. So $\Psi(x,y)\le\sum_{n\ y\text{-smooth}}(\frac xn)^\sigma=x^\sigma\prod_{p\le y}(1-p^{-\sigma})^{-1}$, the sum over all
$y$-smooth $n$ being an Euler product with finitely many factors.

(b) Since $\sigma\ge\frac23$, $\sum_p\sum_{k\ge2}\frac{p^{-k\sigma}}k\le\sum_p\frac{p^{-4/3}}{1-2^{-2/3}}$ is bounded by a constant. Next, with $t=\frac{\log p}{\log y}\in[0,1]$,
$p^{-\sigma}=\frac1pe^{\eta t}\le\frac1p\bigl(1+(e^\eta-1)\frac{\log p}{\log y}\bigr)$, so by Mertens's first and second theorems (Volume~5, Theorems~1.3 and~1.5)
\[
\sum_{p\le y}p^{-\sigma}\le\log\log y+c_1+(e^\eta-1)\frac{\log y+c_2}{\log y}\le\log\log y+c_1+\Bigl(1+\frac{c_2}{\log2}\Bigr)e^\eta.
\]
Exponentiating, $\prod_{p\le y}(1-p^{-\sigma})^{-1}\le C\log y\,e^{Ce^\eta}$, and $x^\sigma=xe^{-\eta\log x/\log y}=xe^{-\eta u}$.

(c) With $\eta=\log u$ (allowed when $u>1$ and $\log u\le\frac13\log y$), $e^{-\eta u+Ce^\eta}=e^{-u\log u+Cu}$. For $y=e^{\sqrt{\log x}}$, $u=\log y=\sqrt{\log x}$, and the condition
$\frac12\log\log x\le\frac13\sqrt{\log x}$ holds for large $x$. Then $\frac{\Psi(x,y)}x\le C\sqrt{\log x}\,e^{-\frac12\sqrt{\log x}\log\log x+C\sqrt{\log x}}=e^{-(\frac12+o(1))\sqrt{\log x}\log\log x}$.

(d) The best bounds miss by factors $\frac{9.95\cdot10^7}{8.80\cdot10^6}=11.3$ and $\frac{1.26\cdot10^{11}}{6.47\cdot10^9}=19.4$. The loss is the price of (a): the smooth $n>x$ are counted with weights
$(\frac xn)^\sigma$ that decay slowly, and a whole range of $n<x$ is counted with weights larger than $1$. The product $\prod_{p\le y}(1-\frac1p)^{-1}\sim e^\gamma\log y$ already shows the
factor $\log y$ that the method cannot remove.

\tools{Volume~5, Theorems~1.3 and~1.5 (Mertens); Euler products (Volume~4).}
\leads Rankin (1938) used this trick to find large gaps between primes. Refined with a saddle point, the same idea gives $\Psi(x,y)$ asymptotically in a huge range (Hildebrand and Tenenbaum):
the loss in (d) is a factor of order $\sqrt u\,\log y$, which the saddle-point method recovers.
\end{solution}

\begin{solution}{pb:dickmanint}
\idea The delay equation $u\rho'(u)=-\rho(u-1)$ becomes an ordinary differential equation for the Laplace transform, because multiplying by $u$ is differentiating in $s$ and the delay is
multiplication by $e^{-s}$. The constant of integration is fixed at $s=\infty$ by $\rho(0)=1$, and Euler's constant appears from the asymptotics of the exponential integral.

\writeup (a) For $u>1$, $\frac{\mathrm d}{\mathrm du}(u\rho(u))=\rho(u)+u\rho'(u)=\rho(u)-\rho(u-1)=\frac{\mathrm d}{\mathrm du}\int_{u-1}^u\rho$, and both sides equal $1$ at $u=1$. If $\rho$ had a zero,
let $u_0>1$ be the first; then $\rho>0$ on $[u_0-1,u_0[$ and $u_0\rho(u_0)=\int_{u_0-1}^{u_0}\rho>0$, a contradiction. So $\rho>0$, and $\rho'(u)=-\frac{\rho(u-1)}u<0$ for $u>1$. For
$u\ge2$, $\rho$ decreases on $[u-1,u]$, so $u\rho(u)\le\rho(u-1)$ and $\rho(u)\le\frac12\rho(u-1)$. In particular $\rho(u)\le2^{2-u}$, so $\hat\rho(s)$ is finite and smooth for
$s>-\log2$, with derivatives obtained under the integral sign.

(b) On $]0,1[$, $u\rho(u)=u$ has derivative $1=\rho(u)-\rho(u-1)$, as $\rho(u-1)=0$; on $]1,\infty[$ this is (a). Since $u\rho(u)$ is continuous and decays exponentially,
\[
\int_0^\infty e^{-su}(u\rho(u))'\dd u=\Bigl[e^{-su}u\rho(u)\Bigr]_0^\infty+s\int_0^\infty ue^{-su}\rho(u)\dd u=-s\hat\rho'(s),
\]
while the left side is $\hat\rho(s)-\int_1^\infty e^{-su}\rho(u-1)\dd u=(1-e^{-s})\hat\rho(s)$.

(c) $\hat\rho>0$, and $(\log\hat\rho)'(s)=-\frac{1-e^{-s}}s$ for $s>0$; the right side extends continuously to $s=0$, and $\hat\rho$ is continuous on $[0,\infty[$, so integrating from $0$ gives the formula.

(d) $s\hat\rho(s)=s\int_0^1e^{-su}\dd u+s\int_1^\infty e^{-su}\rho(u)\dd u=1-e^{-s}+\theta$ with $0\le\theta\le s\int_1^\infty e^{-su}\dd u=e^{-s}$, so $s\hat\rho(s)\to1$. Next,
\[
\int_0^s\frac{1-e^{-t}}t\dd t=\int_0^1\frac{1-e^{-t}}t\dd t+\log s-\int_1^s\frac{e^{-t}}t\dd t\to\log s+\gamma+o(1),
\]
because $\int_0^1\frac{1-e^{-t}}t\dd t-\int_1^\infty\frac{e^{-t}}t\dd t=\gamma$. To see this, integrate by parts in $\Gamma'(1)=\int_0^\infty e^{-t}\log t\dd t=-\gamma$ (Volume~4, Exercise~5.7):
$\int_0^1e^{-t}\log t\dd t=\bigl[(1-e^{-t})\log t\bigr]_0^1-\int_0^1\frac{1-e^{-t}}t\dd t=-\int_0^1\frac{1-e^{-t}}t\dd t$ and $\int_1^\infty e^{-t}\log t\dd t=\int_1^\infty\frac{e^{-t}}t\dd t$. Finally, by (c),
\[
\int_0^\infty\rho=\hat\rho(0)=s\hat\rho(s)\cdot\exp\Bigl(\int_0^s\frac{1-e^{-t}}t\dd t-\log s\Bigr)\to1\cdot e^\gamma.
\]

(e) By Abel summation, $\sum_{n\ y\text{-smooth}}\frac1n=\int_{1^-}^\infty\frac{\mathrm d\Psi(t,y)}t=\int_1^\infty\frac{\Psi(t,y)}{t^2}\dd t$, since $\Psi(t,y)/t\to0$. With $\Psi(t,y)\approx t\rho(\frac{\log t}{\log y})$ and
$t=y^u$, $\frac{\mathrm dt}t=\log y\dd u$, this is $\approx\log y\int_0^\infty\rho(u)\dd u=e^\gamma\log y$, which is Mertens's third theorem. So $\int\rho=e^\gamma$ is the
consistency condition between Dickman's function and Mertens's theorem.

\tools{Volume~4, Exercise~5.7 ($\Gamma'(1)=-\gamma$); Laplace transforms and differentiation under the integral sign (Volume~4); Volume~5, Theorem~1.6.}
\leads The formula of (c) is the starting point of de Bruijn's asymptotic analysis of $\rho$, which gives $\rho(u)=e^{-u\log u-u\log\log u+u+o(u)}$ by a saddle point applied to the inverse Laplace
transform. Its counterpart for permutations is the probability that a random permutation has no cycle longer than $\frac nu$, which also tends to $\rho(u)$.
\end{solution}

\begin{solution}{pb:buchstab}
\idea A rough number below $x$ with no prime factor below $x^{1/3}$ is $1$, a prime, or a product of two primes; counting the products with the prime number theorem gives
$\int\frac{\mathrm dt}{t(1-t)}$. Independent sieving would give $e^{-\gamma}$ everywhere, and the actual answer is far from it when $y$ is a large power of $x$.

\writeup (a) A $y$-rough $n\le x$ with $n>1$ and $y\ge x^{1/3}$ has at most two prime factors, since three would multiply to more than $y^3\ge x$. If $y\ge\sqrt x$ it cannot have two
either, so it is a prime in $]y,x]$. If $x^{1/3}\le y<\sqrt x$, it is a prime in $]y,x]$ or a product $pq$ with $y<p\le q$, $pq\le x$.

(b) For $1<u\le2$, $\Phi(x,y)=1+\pi(x)-\pi(y)$, which is $\sim\frac x{\log x}$ since $\pi(x^{1/u})=o(\pi(x))$. Multiplying by $\frac{\log y}x=\frac{\log x}{ux}$ gives $\frac1u$.
For $2<u\le3$, the pairs number
\[
\sum_{y<p\le\sqrt x}\Bigl(\pi\Bigl(\frac xp\Bigr)-\pi(p)+1\Bigr),
\]
which is $\sum_{y<p\le\sqrt x}\pi(\frac xp)+O(\pi(\sqrt x)^2)$.
Since $\frac xp\ge\sqrt x$, $\pi(\frac xp)=(1+o(1))\frac{x}{p\log(x/p)}$ uniformly, and the sum is
\[
(1+o(1))\frac x{\log x}\sum_{y<p\le\sqrt x}\frac1pF\Bigl(\frac{\log p}{\log x}\Bigr)\qquad\text{with }F(t)=\frac1{1-t}.
\] By Mertens's
second theorem (Volume~5, Theorem~1.5), $\sum_{x^a<p\le x^b}\frac1p\to\log\frac ba$ for fixed $0<a<b$; cutting $[\frac1u,\frac12]$ into small intervals and using the continuity of $F$, Riemann sums give
\[
\sum_{x^{1/u}<p\le x^{1/2}}\frac1pF\Bigl(\frac{\log p}{\log x}\Bigr)\to\int_{1/u}^{1/2}\frac{\mathrm dt}{t(1-t)}=\Bigl[\log\frac t{1-t}\Bigr]_{1/u}^{1/2}=\log(u-1).
\]
So $\Phi(x,y)=(1+\log(u-1)+o(1))\frac x{\log x}$ and $\frac{\log y}x\Phi(x,y)\to\frac{1+\log(u-1)}u$. (At $u=2$ both formulas give $\frac12$.)

(c) For $2<u<3$, $(u\omega(u))'=\frac1{u-1}=\omega(u-1)$, as $u-1\in\left]1,2\right[$. Numerically $\omega(2)=0.5<e^{-\gamma}=0.56146<\omega(3)=\frac{1+\log2}3=0.56438$.

(d) $\omega$ decreases from $1$ to $\frac12$ on $]1,2]$; on $[2,3]$, $\omega(u)\ge\frac12$ because $2+2\log(u-1)-u$ vanishes at $u=2$ and has derivative $\frac2{u-1}-1\ge0$. So the minimum on $]1,3]$ is
$\omega(2)=\frac12$, where independent sieving overestimates by the factor $2e^{-\gamma}=1.1229$. This is Legendre's failure: sieving the integers up to $x$ by the primes up to $\sqrt x$
leaves the primes of $]\sqrt x,x]$, about $\frac x{\log x}$ of them, not $x\prod_{p\le\sqrt x}(1-\frac1p)\sim2e^{-\gamma}\frac x{\log x}$ (Volume~5, Exercise~8.2). As $u\to1$, on the other hand, the prediction
underestimates by a factor tending to $e^\gamma$.

\tools{Volume~5, Theorems~1.5 and~5.5, Exercise~8.2.}
\leads Buchstab (1937) found the equation $(u\omega(u))'=\omega(u-1)$, and de Bruijn (1950) proved that $\omega(u)\to e^{-\gamma}$, the difference oscillating and decaying faster than any exponential
(Figure~\ref{fig:buchstab}). So independent sieving is right exactly when the primes used are a small power of the range: for $y=x^{1/u}$ with $u$ large. This is the ``sieve limit'' of Volume~5,
Chapter~8, seen from the exact side.
\end{solution}

\begin{solution}{pb:polymertens}
\idea In $\mathbb F_q[x]$ the prime number theorem is an identity, $\sum_{k\mid m}kN_k=q^m$, so the logarithm of Mertens's product is \emph{exactly} the harmonic sum, up to a tiny
error. Euler's constant enters through $H_n-\log n\to\gamma$, the polynomial counterpart of the integral $\int_0^1\frac{1-e^{-t}}t-\int_1^\infty\frac{e^{-t}}t$ of Problem~\ref{pb:dickmanint}.

\writeup (a) $-\log E_n=\sum_{k\le n}N_k\sum_{j\ge1}\frac{q^{-jk}}j$. The terms with $jk=m\le n$ give, for each $m$, $\frac{q^{-m}}m\sum_{k\mid m}kN_k=\frac1m$; together $H_n$.
The others have $k\le n<jk$, so $j\ge2$. Using $N_k\le\frac{q^k}k$, each is at most $\frac{q^{-(j-1)k}}{jk}$. For fixed $k$ the smallest admissible $j$, say $j_0$, has $(j_0-1)k\ge\frac n2$:
if $k\le\frac n2$ then $(j_0-1)k=\lfloor\frac nk\rfloor k>n-k\ge\frac n2$, and if $k>\frac n2$ then $j_0=2$ and $(j_0-1)k=k>\frac n2$. The series over $j$ has ratio at most $q^{-k}\le\frac12$, so
\[
0\le\varepsilon_n\le\sum_{k\le n}\frac{2q^{-n/2}}{j_0k}\le q^{-n/2}H_n\le(1+\log n)q^{-n/2}.
\]
Hence $E_n=e^{-H_n-\varepsilon_n}\sim e^{-\log n-\gamma}=\frac{e^{-\gamma}}n$.

(b) A reducible polynomial of degree $n$ has a factor of degree at most $\frac n2$, so the polynomials with no such factor are the irreducible ones, a proportion $\frac1n+O(\frac{q^{-n/2}}n)$
(Problem~\ref{pb:randpoly}). The prediction is $E_{\lfloor n/2\rfloor}\sim\frac{e^{-\gamma}}{n/2}=\frac{2e^{-\gamma}}n$. The ratio tends to $2e^{-\gamma}$.

(c) $H_n=\log n+\gamma+\frac1{2n}-\frac1{12n^2}+O(\frac1{n^4})$ (Euler--Maclaurin, Volume~4, Corollary~11.6), so $ne^\gamma E_n=\exp\bigl(-\frac1{2n}+\frac1{12n^2}+O(\frac1{n^4})-\varepsilon_n\bigr)=1-\frac1{2n}+O(\frac1{n^2})$.
At $n=10$, $\exp(-\frac1{20}+\frac1{1200})=0.952017$; for $q=5$, $\varepsilon_{10}\le(1+\log10)5^{-5}$ is small, and the numerical value $0.952015$ agrees.

\tools{Counting irreducible polynomials (Volume~2, Chapter~9); Volume~4, Corollary~11.6; Problem~\ref{pb:randpoly}.}
\leads The same factor $2e^{-\gamma}$ in both worlds shows that Legendre's failure is not a defect of the integers: it is what happens whenever one sieves an interval with ``primes'' up to the
square root of its size. The analogy between $\mathbb Z$ and $\mathbb F_q[x]$, with $\log x$ playing the role of the degree, runs through the whole of number theory; Rosen's book develops it.
\end{solution}

\begin{solution}{pb:hawkins}
\idea Hawkins's sieve is Eratosthenes's sieve with the arithmetic replaced by coin tosses: the strikes on different integers are independent. Then the survival probability of $n$ is
literally the sieve product, and a one-line recurrence shows that it behaves like $\frac1{\log n}$.

\writeup (a) The past determines which $X_k$ are below $n$. The strikes aimed at $n$, one for each $X_k<n$, are independent of the past and of each other, and each spares $n$ with probability
$1-\frac1{X_k}$. If $n$ survives them all, it is the least survivor above the previous $X_k$, so it is a Hawkins prime. Hence the conditional probability is $Q_n$. Then $Q_{n+1}=Q_n(1-\frac1n)$
if $n$ is a Hawkins prime and $Q_{n+1}=Q_n$ otherwise, so $\mathbb E[Q_{n+1}\mid\text{past}]=Q_n-\frac{Q_n}n\cdot Q_n$.

(b) Taking expectations and using $\mathbb E[Q_n^2]\ge\mathbb E[Q_n]^2$, $e_n=\mathbb E[Q_n]$ satisfies $e_{n+1}\le e_n-\frac{e_n^2}n=f_n(e_n)$, with $f_n(t)=t-\frac{t^2}n$ increasing on
$[0,1]$ for $n\ge2$. Since $e_2=1=q_2$ (there is no $X_k<2$), induction gives $e_{n+1}\le f_n(e_n)\le f_n(q_n)=q_{n+1}$, provided $0\le q_n\le1$, which also follows by induction.

(c) $\frac1{q_n(1-q_n/n)}=\frac n{q_n(n-q_n)}=\frac1{q_n}+\frac1{n-q_n}$. As $0<q_k\le1$, $\frac1k\le\frac1{k-q_k}\le\frac1{k-1}$, so $\frac1{q_n}=1+\sum_{k=2}^{n-1}\frac1{k-q_k}$ lies
between $H_{n-1}$ and $1+H_{n-2}$: $\frac1{q_n}=\log n+O(1)$. The expected number of Hawkins primes up to $x$ is $\sum_{n\le x}e_n\le\sum_{n\le x}q_n=\sum_{n\le x}\frac1{\log n+O(1)}\sim\frac x{\log x}$.
(Numerically, one finds $\frac1{q_n}-\log n\to1.2400$.)

(d) The strikes of Eratosthenes's sieve are not independent: $n$ is struck by $p$ exactly when $p\mid n$, and the events $p\mid n$ for different $p$ are independent only on average over
intervals much longer than $\prod p$, by the Chinese remainder theorem. For $p\le\sqrt x$ that product is $e^{(1+o(1))\sqrt x}$, astronomically longer than $[1,x]$. Over such a short
interval the sieve leaves fewer survivors than the product predicts, by the factor $2e^{-\gamma}$ of Problem~\ref{pb:buchstab}. Hawkins's coins know nothing of the Chinese remainder
theorem, so for them the product is the truth.

\tools{Conditional expectation and Jensen's inequality (Volume~3, Theorem~4.14); the harmonic numbers (Volume~0).}
\leads With probability $1$, the Hawkins primes satisfy analogues of the prime number theorem and of Mertens's theorem (Wunderlich, 1974 and 1976), and even of the Riemann
hypothesis (Neudecker and Williams, 1974; Heyde, 1976, with martingales). Random models
like this one and Cram\'er's (Volume~5, Chapter~10) are good at the large-scale count of primes and wrong about arithmetic: Hawkins primes are as often even as odd.
\end{solution}

\begin{solution}{pb:renyi}
\idea $z^{\Omega(n)-\omega(n)}$ is a multiplicative function that is $1$ on primes and differs from $1$ only at prime powers, so it is a convolution $\mathbf1*h$ with $h$ supported on
squarefull numbers, and Wintner's theorem gives its mean as an Euler product. The pole of the factor at $p=2$ decides the tail of the law.

\writeup (a) $\Omega-\omega$ is additive, $\Omega(p^k)-\omega(p^k)=k-1$, so $f_z$ is multiplicative with $f_z(p^k)=z^{k-1}$. Then $h_z=f_z*\mu$ is multiplicative with
$h_z(p)=f_z(p)-1=0$ and $h_z(p^k)=z^{k-1}-z^{k-2}=z^{k-2}(z-1)$ for $k\ge2$. For $|z|<2$,
\[
\sum_d\frac{|h_z(d)|}d=\prod_p\Bigl(1+|z-1|\sum_{k\ge2}\frac{|z|^{k-2}}{p^k}\Bigr)=\prod_p\Bigl(1+\frac{|z-1|}{p(p-|z|)}\Bigr)<\infty,
\]
since $p-|z|>0$ for all $p\ge2$ and $\sum_p\frac1{p^2}$ converges.

(b) By Wintner's theorem the mean value is $\sum_d\frac{h_z(d)}d=\prod_p\bigl(1+\frac{z-1}{p(p-z)}\bigr)$. And
$(1-\frac1p)(1+\frac1{p-z})=\frac{(p-1)(p-z+1)}{p(p-z)}=\frac{p^2-pz+z-1}{p(p-z)}=1+\frac{z-1}{p(p-z)}$.

(c) Let $d_k(x)=\frac1x\#\{n\le x:\Omega(n)-\omega(n)=k\}$ and $P_x(z)=\sum_kd_k(x)z^k=\frac1x\sum_{n\le x}z^{\Omega(n)-\omega(n)}$, a polynomial. By Cauchy's formula on the unit circle,
$d_k(x)=\frac1{2\pi}\int_0^{2\pi}P_x(e^{i\theta})e^{-ik\theta}\dd\theta$. By (b) $P_x(e^{i\theta})\to R(e^{i\theta})$ for every $\theta$, and $|P_x|\le\frac{\lfloor x\rfloor}x\le1$ there, so by dominated
convergence $d_k(x)\to\frac1{2\pi}\int R(e^{i\theta})e^{-ik\theta}\dd\theta=[z^k]R$, $R$ being holomorphic on $|z|<2$ (Volume~4, Theorem~9.10). So $d_k$ exists and $\sum d_kz^k=R(z)$.
Then $d_0=R(0)=\prod_p(1-\frac1{p^2})=\frac1{\zeta(2)}=\frac6{\pi^2}$: the squarefree numbers, as it should be. The mean is $R'(1)=\sum_p\frac{\mathrm d}{\mathrm dz}\log\bigl(1+\frac{z-1}{p(p-z)}\bigr)\big|_{z=1}=\sum_p\frac1{p(p-1)}$,
since $R(1)=1$ and the derivative of $\frac{z-1}{p(p-z)}$ at $1$ is $\frac1{p(p-1)}$.

(d) Write $R(z)=\frac12\bigl(1+\frac1{2-z}\bigr)R_3(z)$ with $R_3(z)=\prod_{p\ge3}(1-\frac1p)(1+\frac1{p-z})$, holomorphic on $|z|<3$. Then
$R(z)-\frac{R_3(2)}{2(2-z)}$ is holomorphic on $|z|<3$, so by Cauchy's estimates (Volume~4, Corollary~9.12) its coefficients are $O(2.5^{-k})$, while $\frac{R_3(2)}{2(2-z)}=\frac{R_3(2)}4\sum_k\frac{z^k}{2^k}$. Hence
$2^kd_k\to\frac{R_3(2)}4$, and
\[
R_3(2)=\prod_{p\ge3}\frac{p-1}p\cdot\frac{p-1}{p-2}=\prod_{p\ge3}\frac{(p-1)^2}{p(p-2)}=\prod_{p\ge3}\Bigl(1-\frac1{(p-1)^2}\Bigr)^{-1}=\frac1{C_2}.
\]

(e) The agreement is to five or six decimals: the convergence of $\frac1x\sum_{n\le x}f(n)$ to the mean is fast here, with error $O(x^{-1/2+\varepsilon})$ since $h$ lives on squarefull numbers.

\tools{Volume~5, Theorem~2.2 (Wintner); Volume~4, Theorem~9.10 and Corollary~9.12 (Cauchy's formula and estimates).}
\leads R\'enyi's law is a law of repeated prime factors, governed by the small primes, unlike the Erd\H os--Kac law of $\omega(n)$, governed by all of them. The twin prime constant enters
through the factor $\frac{p-1}{p-2}$, the same factor that makes pairs $p,p+2$ more common than independent coins would (Problem~\ref{pb:digitpairs}).
\end{solution}

\begin{solution}{pb:inspection}
\idea A random integer is more likely to fall in a long gap than in a short one, in proportion to the length. So the gap seen from a random point has mean $\frac{\sum g^2}{\sum g}$,
and Cauchy--Schwarz already makes this at least the average gap.

\writeup (a) Let $N=\pi(x)$. The gaps with $p_{n+1}\le x$ are $g_1,\dots,g_{N-1}$, with sum $p_N-2$. By Cauchy--Schwarz,
$\sum_{n<N}g_n^2\ge\frac{(p_N-2)^2}{N-1}$. By the prime number theorem $N\sim\frac x{\log x}$, and $p_N\sim x$ (since $p_{N+1}>x$ and $p_{N+1}\sim p_N$). Hence $\sum g_n^2\ge(1+o(1))x\log x$.

(b) Each gap $]p_n,p_{n+1}]$ contains $g_n$ integers, each of which sees $G=g_n$. So $\sum_{2<m\le p_N}G(m)=\sum_{n<N}g_n^2$, and dividing by $p_N-2$ gives $\mathbb E[G]$; by (a) it is at least
$(1+o(1))\frac{x\log x}x=(1+o(1))\log x$. The average gap $\frac{p_N-2}{N-1}$ is $\sim\log x$: the gap containing a random point is size-biased, as the bus interval containing a random
arrival time in Volume~3, Proposition~9.10.

(c) With $P(G_0=k)=q(1-q)^{k-1}$, $\mathbb E[G_0]=\frac1q$ and $\mathbb E[G_0(G_0-1)]=\frac{2(1-q)}{q^2}$, so $\mathbb E[G_0^2]=\frac{2-q}{q^2}$ and the ratio is $\frac{2-q}q$. With $q=\frac1{\log x}$,
the size-biased mean is $\sim2\log x$, and $\sum g_n^2\approx(\text{number of gaps})\cdot\mathbb E[G_0^2]\approx\frac x{\log x}\cdot2\log^2x=2x\log x$.

(d) The ratio climbs slowly, $1.16$, $1.52$, $1.63$, well above the proved bound $1$ and still well below $2$. It is consistent with the prediction $2$ approached with corrections of order
$\frac1{\log x}$ (the gaps near $10^8$ are not yet exponential: gaps of $2$ and $4$ are over-represented and the mean gap $\log x$ is only $18$). If one draws $m$ at random and takes the next prime,
the primes that follow long gaps are favoured, with probability proportional to the preceding gap: about twice the uniform weight on average, if Cram\'er is right.

\tools{Volume~3, Proposition~2.6 (Cauchy--Schwarz) and Proposition~9.10 (inspection paradox); Volume~5, Theorem~5.5 and Chapter~10.}
\leads Some prime generators used in cryptography draw a random number and search upwards for the next prime (Brandt and Damg\aa rd, 1992). By (b) their output is not uniform among
primes; Fouque and Tibouchi (2014) showed, for instance, that the upper member of a twin pair is output much less often than a uniform choice would. Proving
$\sum_{p_{n+1}\le x}g_n^2=O(x^{1+\varepsilon})$ without the Riemann hypothesis is an open problem.
\end{solution}

\begin{solution}{pb:digitpairs}
\idea In the model the next ``prime'' arrives after a geometric number of steps along the cycle $1\to3\to7\to9$, so the nearest endings are favoured and a repeat, needing a full turn,
is the least likely. This explains the deficit of repeats. It cannot explain the excess of $1\to7$, which comes from divisibility by $3$, invisible to a model that only knows
about $2$ and $5$.

\writeup (a) The next ``prime'' is at step $j+4t$ along the cycle with probability $q(1-q)^{j+4t-1}$, for $t\ge0$ and $j\in\{1,2,3,4\}$, and its ending is determined by $j$.
Summing over $t$, $\pi_j=q(1-q)^{j-1}\sum_t(1-q)^{4t}=\frac{q(1-q)^{j-1}}{1-(1-q)^4}$.

(b) $1-(1-q)^4=4q-6q^2+O(q^3)=4q(1-\frac32q+O(q^2))$ and $q(1-q)^{j-1}=q(1-(j-1)q+O(q^2))$, so
$\pi_j=\frac14\bigl(1-(j-1)q+\frac32q+O(q^2)\bigr)=\frac14+\frac{(5-2j)q}8+O(q^2)$. The rarest is $j=4$, the repeat: $\pi_4=\frac14-\frac{3q}8+O(q^2)$.

(c) Near $x$ the density of primes is $\frac1{\log x}$ and the integers prime to $10$ have density $\frac4{10}$, so among them the primes have density $\frac{10}{4\log x}$. With $q=\frac{2.5}{\log p}$ the
model predicts a repeat probability $19.5\%$ instead of $25\%$; the primes show $17.7\%$. So the cyclic order alone accounts for $5.5$ of the $7.3$ percentage points of deficit, about three quarters.

(d) The transitions $1\to3$, $1\to7$, $1\to9$ need gaps $\equiv2,6,8\pmod{10}$, the smallest being $2$, $6$, $8$. By Hardy and Littlewood, prime pairs at distance $h$ have density
proportional to $\prod_{p\mid h,\,p>2}\frac{p-1}{p-2}$, which is $1$ for $h=2,4,8$ and $2$ for $h=6$: since $p$ and $p+6$ are automatically in the same class modulo~$3$, the condition
``neither is divisible by $3$'' costs a factor $\frac23$ instead of $\frac13$. The coin model, blind to the prime $3$, misses the doubled weight of gaps divisible by $6$, and so undercounts $1\to7$.

\tools{Geometric series (Volume~0); Volume~5, Chapter~10 (Cram\'er's model and the Hardy--Littlewood constants).}
\leads Lemke Oliver and Soundararajan (2016) derived from the Hardy--Littlewood conjectures a precise prediction: each transition frequency is $\frac14$ plus a term of size
$\frac{\log\log x}{\log x}$ plus a term of size $\frac1{\log x}$, so the bias does disappear, but only at the speed of $\frac1{\log x}$, which is why it is so visible at $10^8$.
\end{solution}

\section*{Map of the chapter}

\begin{center}
\small
\adjustbox{max width=\linewidth}{\begin{tabular}{lll}
Problem & Topic & Volumes used\\\hline
\ref{pb:leastpf} & sums of primes and of least prime factors & 5\\
\ref{pb:largestpf} & the largest prime factor on average & 4, 5\\
\ref{pb:radical} & the average radical & 4, 5\\
\ref{pb:factorialdiv} & the number of divisors of $n!$ & 0, 1, 5\\
\ref{pb:centralbinom} & the prime factors of $\binom{2n}n$ & 1, 5\\
\ref{pb:primezeta} & $\sum_pp^{-2}$ by Möbius inversion & 4, 5\\
\ref{pb:superprimes} & primes with a prime index & 5\\
\ref{pb:ramprimes} & Ramanujan primes & 1, 5\\
\ref{pb:pminusone} & squarefree $p-1$ and Artin's constant & 1, 4, 5\\
\ref{pb:titchmarsh} & the divisors of $p-1$ & 5\\
\ref{pb:menon} & roots of $x^n-1$ modulo $p$, Menon's identity & 1, 2, 5\\
\ref{pb:omegapoly} & the prime factors of $n^2+1$ & 1, 3, 5\\
\ref{pb:rankin} & Rankin's bound for smooth numbers & 4, 5\\
\ref{pb:dickmanint} & $\int_0^\infty\rho=e^\gamma$ & 4, 5\\
\ref{pb:buchstab} & rough numbers, Buchstab's function & 5\\
\ref{pb:polymertens} & Mertens's theorem in $\mathbb F_q[x]$ & 2, 4, 5\\
\ref{pb:hawkins} & Hawkins's random sieve & 0, 3, 5\\
\ref{pb:renyi} & the law of $\Omega(n)-\omega(n)$ & 4, 5\\
\ref{pb:inspection} & the inspection paradox for prime gaps & 3, 5\\
\ref{pb:digitpairs} & last digits of consecutive primes & 0, 5\\
\end{tabular}}
\end{center}

\section*{Sources and further reading}

\begin{itemize}
\item D.~Zagier, The first 50 million prime numbers, \emph{The Mathematical Intelligencer} 0 (1977), 7--19, translated by R.~Perlis from the author's inaugural lecture at Bonn, 5~May 1975 (the epigraph).
\item R.~J.~Lemke Oliver and K.~Soundararajan, Unexpected biases in the distribution of consecutive primes, \emph{Proceedings of the National Academy of Sciences} 113 (2016), E4446--E4454;
E.~Klarreich, Mathematicians discover prime conspiracy, \emph{Quanta Magazine}, 13~March 2016 (Tokieda's lecture and the quotation of Lemke Oliver).
\item K.~Alladi and P.~Erd\H os, On an additive arithmetic function, \emph{Pacific Journal of Mathematics} 71 (1977), 275--294.
\item D.~W.~Masser, Open problems, in W.~W.~L.~Chen (ed.), \emph{Proceedings of the Symposium on Analytic Number Theory}, Imperial College, London, 1985; J.~Oesterl\'e, Nouvelles approches du
``th\'eor\`eme'' de Fermat, S\'eminaire Bourbaki 694, \emph{Ast\'erisque} 161--162 (1988), 165--186.
\item P.~Erd\H os, S.~W.~Graham, A.~Ivi\'c and C.~Pomerance, On the number of divisors of $n!$, in \emph{Analytic Number Theory, Vol.~1} (Allerton Park, 1995), Progress in Mathematics 138,
Birkh\"auser, 1996, 337--355.
\item P.~Erd\H os, Beweis eines Satzes von Tschebyschef, \emph{Acta Scientiarum Mathematicarum (Szeged)} 5 (1932), 194--198.
\item S.~Ramanujan, A proof of Bertrand's postulate, \emph{Journal of the Indian Mathematical Society} 11 (1919), 181--182; J.~Sondow, Ramanujan primes and Bertrand's postulate, \emph{American
Mathematical Monthly} 116 (2009), 630--635; S.~Laishram, On a conjecture on Ramanujan primes, \emph{International Journal of Number Theory} 6 (2010), 1869--1873.
\item R.~E.~Dressler and S.~T.~Parker, Primes with a prime subscript, \emph{Journal of the ACM} 22 (1975), 380--381.
\item C.~Hooley, On Artin's conjecture, \emph{Journal f\"ur die reine und angewandte Mathematik} 225 (1967), 209--220; P.~J.~Stephens, An average result for Artin's conjecture, \emph{Mathematika} 16 (1969),
178--188; P.~Moree, Artin's primitive root conjecture, a survey, \emph{Integers} 12A (2012), A13.
\item E.~C.~Titchmarsh, A divisor problem, \emph{Rendiconti del Circolo Matematico di Palermo} 54 (1930), 414--429; Ju.~V.~Linnik, \emph{The Dispersion Method in Binary Additive Problems},
American Mathematical Society, 1963; H.~Halberstam, Footnote to the Titchmarsh--Linnik divisor problem, \emph{Proceedings of the AMS} 18 (1967), 187--188.
\item P.~K.~Menon, On the sum $\sum(a-1,n)$ [$(a,n)=1$], \emph{Journal of the Indian Mathematical Society} 29 (1965), 155--163; P.~Stevenhagen and H.~W.~Lenstra, Chebotar\"ev and his density theorem,
\emph{The Mathematical Intelligencer} 18:2 (1996), 26--37.
\item R.~A.~Rankin, The difference between consecutive prime numbers, \emph{Journal of the London Mathematical Society} 13 (1938), 242--247; A.~Hildebrand and G.~Tenenbaum, Integers without large prime
factors, \emph{Journal de Th\'eorie des Nombres de Bordeaux} 5 (1993), 411--484.
\item K.~Dickman, On the frequency of numbers containing prime factors of a certain relative magnitude, \emph{Arkiv f\"or Matematik, Astronomi och Fysik} 22A (1930), 1--14; N.~G.~de Bruijn,
On the number of positive integers $\le x$ and free of prime factors $>y$, \emph{Proceedings of the Koninklijke Nederlandse Akademie van Wetenschappen} 54 (1951), 50--60;
G.~Tenenbaum, \emph{Introduction to Analytic and Probabilistic Number Theory}, 3rd edition, American Mathematical Society, 2015, Chapter~III.5--6.
\item A.~A.~Buchstab, Asymptotic estimates of a general number-theoretic function, \emph{Matemati\v cheski\u\i\ Sbornik} 2(44) (1937), 1239--1246; N.~G.~de Bruijn, On the number of uncancelled elements in
the sieve of Eratosthenes, \emph{Proceedings of the Koninklijke Nederlandse Akademie van Wetenschappen} 53 (1950), 803--812.
\item M.~Rosen, \emph{Number Theory in Function Fields}, Graduate Texts in Mathematics 210, Springer, 2002.
\item D.~Hawkins, The random sieve, \emph{Mathematics Magazine} 31 (1957/58), 1--3; M.~C.~Wunderlich, A probabilistic setting for prime number theory, \emph{Acta Arithmetica} 26 (1974), 59--81, and
The prime number theorem for random sequences, \emph{Journal of Number Theory} 8 (1976), 369--371; W.~Neudecker and D.~Williams, The ``Riemann hypothesis'' for the Hawkins random sieve,
\emph{Compositio Mathematica} 29 (1974), 197--200; C.~C.~Heyde, On asymptotic behavior for the Hawkins random sieve, \emph{Proceedings of the AMS} 56 (1976), 277--280.
\item A.~R\'enyi, On the density of certain sequences of integers, \emph{Publications de l'Institut Math\'ematique (Belgrade)} 8 (1955), 157--162.
\item H.~Cram\'er, On the order of magnitude of the difference between consecutive prime numbers, \emph{Acta Arithmetica} 2 (1936), 23--46; G.~H.~Hardy and J.~E.~Littlewood, Some problems of
``Partitio numerorum'' III, \emph{Acta Mathematica} 44 (1923), 1--70.
\item J.~Brandt and I.~Damg\aa rd, On generation of probable primes by incremental search, in \emph{Advances in Cryptology, CRYPTO '92}, Lecture Notes in Computer Science 740, Springer, 1993,
358--370; P.-A.~Fouque and M.~Tibouchi, Close to uniform prime number generation with fewer random bits, in \emph{Automata, Languages, and Programming (ICALP 2014)}, Springer, 2014.
\end{itemize}
